% =================================================
\section{Requisitos e Regras de Negócio}
% =================================================

\subsection{Requisitos funcionais}


\setlength{\tabcolsep}{5pt}
\begin{longtable}{p{0.10\textwidth} p{0.82\textwidth}}
\toprule
\textbf{ID} & \textbf{Requisito}\\
\midrule
\endfirsthead
\toprule
\textbf{ID} & \textbf{Requisito} (continuação)\\
\midrule
\endhead
\bottomrule
\endfoot
RF01 & O sistema deve permitir login por e-mail e senha.\\
RF02 & O sistema deve permitir login com Google através do provedor de autenticação configurado.\\
RF03 & O sistema deve permitir recuperação de senha.\\
RF04 & O sistema deve permitir que um usuário possua múltiplos papéis.\\
RF05 & O sistema deve permitir alternância de papel ativo no dashboard quando houver mais de um papel autorizado.\\
RF06 & O sistema deve permitir cadastro e manutenção de bens patrimoniais conforme permissões.\\
RF07 & O sistema deve manter histórico auditável de alterações patrimoniais.\\
RF08 & O sistema deve permitir requisição de adição de patrimônio.\\
RF09 & O sistema deve permitir requisição de edição de patrimônio.\\
RF10 & O sistema deve impedir mais de uma requisição de edição pendente para o mesmo bem.\\
RF11 & O sistema deve permitir aprovação ou rejeição de requisições.\\
RF12 & O sistema deve permitir registro da baixa de um bem após o procedimento institucional.\\
RF13 & O sistema deve permitir cadastro de almoxarifados.\\
RF14 & O sistema deve calcular pesos e volumes de frascos de reagentes utilizando a densidade.\\
RF15 & O sistema deve registrar empréstimo, devolução, descarte, esvaziamento e quebra de frascos.\\
RF16 & O sistema deve permitir consulta por nome, número de patrimônio, status, local e demais filtros aplicáveis.\\
RF17 & O sistema deve permitir criação de turmas com capacidade máxima estabelecida.\\
RF18 & O sistema deve permitir entrada de alunos em turmas por código ou e-mail.\\
RF19 & O sistema deve permitir criação de posts em turmas pelo professor responsável.\\
RF20 & O sistema deve permitir comentários em posts conforme regras de acesso.\\
RF21 & O sistema deve permitir upload de roteiros em PDF.\\
RF22 & O sistema deve permitir compartilhamento de roteiros entre professores.\\
RF23 & O sistema deve permitir vínculo de roteiro a turma por meio de post.\\
RF24 & O sistema deve gerar relatórios mensais e por período, incluindo consumo em mL e massa em g.\\
RF25 & O sistema deve manter dados históricos e de auditoria sem apagar fatos relevantes.\\
\bottomrule
\end{longtable}

\subsection{Regras de Negócio Gerais}
\subsubsection{Status de um frasco de reagente}
\begin{regra}
Quando um frasco de reagente é cadastrado como em quarentena, deve ser obrigatório escrever o detalhe do status explicando o motivo dele estar em quarentena. Essa justificativa humana obrigatória deve possuir no mínimo 20 caracteres após trim; mensagens automáticas do servidor não são entradas humanas.
Um frasco em quarentena não pode ser descartado diretamente. A quarentena deve ser encerrada por decisão humana explícita: \texttt{VOLTAR\_A\_DISPONIVEL} ou \texttt{PENDENTE\_DE\_DESCARTE}. No segundo caso, a decisão gera autorização operacional estruturada, mantém o frasco \texttt{INDISPONIVEL} e somente então permite \texttt{descartarFrasco}, conforme a máquina de estados da Seção~4.
\end{regra}

% M9 -- fonte normativa única de autorização
\subsection{Autorização e usuários (M9)}
\label{sec:m9-autorizacao}

Esta subseção é a fonte normativa única de autorização. Se uma tabela de
stakeholder, tela, fluxo, pseudocódigo ou Rules divergir, esta subseção
prevalece. Autenticação responde quem apresentou uma credencial Firebase
válida; autorização responde \texttt{podeExecutar(uid, operacao, recurso,
estadoAtual)}. A primeira não implica a segunda.

\paragraph{Modelo e fontes de verdade.}
O modelo é RBAC + escopo/vínculo persistido + ownership quando o domínio o
exige. Os únicos papéis são Chefe\_Geral, Gestor\_Almoxarifado,
Gestor\_Bens\_Patrimoniais, Professor, Aluno e Bolsista. Papéis atuais são
os documentos de papel cujo \texttt{id\_usuario} coincide com o UID; Bolsista
implica Aluno; Chefe\_Geral é exclusivo; Professor e Aluno não coexistem; e
Bolsista e Gestor\_Almoxarifado não coexistem. Conta, papel, vínculo,
ownership, recurso ou escopo ausente, inválido, inconsistente ou desconhecido
significa negação. Nenhum payload, cache, espelho, campo denormalizado, papel
selecionado na UI ou claim antiga cria autorização.

\begin{itemize}
  \item \texttt{Usuarios/uid.ativo=true} é pré-condição universal. A ausência
  do documento ou \texttt{ativo=false} nega.
  \item A relação \texttt{Gestor\_Almoxarifado\_x\_Almoxarifado} é o escopo
  N:N canônico do gestor; um gestor pode ter zero ou muitos vínculos.
  Chefe\_Geral tem escopo global já decidido; Gestor\_Bens\_Patrimoniais tem
  escopo global do domínio patrimonial. Não há escopo patrimonial alternativo
  nesta V1.
  \item Turma pertence a \texttt{id\_professor}; matrícula canônica é
  \texttt{Turma/id/Alunos/uid}; o espelho em
  \texttt{Usuarios/uid/Turmas/id} é somente consulta. Roteiro pertence ao
  professor que o enviou; compartilhamento persistido autoriza apenas o uso
  previsto, não redistribuição.
  \item A decisão de domínio continua separada: RBAC não torna válida
  transição M5 proibida, não corrige fato metrológico M6, não substitui
  idempotência M7 e não contorna cache/estoque/notificação M8.
\end{itemize}

\paragraph{Matriz normativa de autorização.}
Todas as linhas exigem identidade autenticada e Usuário ativo. Tx exige leitura
da autorização e do recurso na mesma transação que produz a escrita. Rules
significa leitura direta de cliente apenas se confirmar versão atual e relação
persistida; backend é Cloud Function privilegiada com validação server-side.

\scriptsize
\setlength{\tabcolsep}{2pt}
\begin{longtable}{>{\raggedright\arraybackslash}p{0.15\textwidth}>{\raggedright\arraybackslash}p{0.18\textwidth}>{\raggedright\arraybackslash}p{0.19\textwidth}>{\raggedright\arraybackslash}p{0.16\textwidth}>{\raggedright\arraybackslash}p{0.17\textwidth}}
\toprule
Ator/papel & Operação/recurso & Escopo ou ownership & Decisão/origem & Resultado \\
\midrule
Chefe\_Geral & gerir usuário, papel, vínculo e almoxarifado & global; RN-ROLE & backend, Tx & permite ou nega invariantes \\
Chefe\_Geral & operar domínios existentes & global, sem dispensar domínio & backend, Tx em mutação & permite somente matriz \\
Gestor\_Almoxarifado & cadastro, retirada, devolução, M5/M6, relatório e etiqueta & vínculo atual com almox do recurso & backend, Tx & nega outro almox \\
Gestor\_Almoxarifado & catálogo global de reagentes & papel atual; sem escopo de almox & backend & permite ações de catálogo V1 \\
Gestor\_Bens\_Patrimoniais & bem, resposta, baixa e relatório patrimonial & domínio patrimonial global & backend, Tx em mutação & permite; professor não substitui gestor \\
Professor & turma, post, roteiro e requisição próprios & professor dono; roteiro próprio/compartilhado & backend, Tx em mutação & nega recurso alheio \\
Aluno & ingresso, leitura/comentário acadêmico & matrícula atual ou convite válido & backend/Rules & nega turma, post ou roteiro alheio \\
Bolsista & direitos de Aluno e consulta/recebimento de empréstimo & matrícula e próprio empréstimo & backend/Rules & não opera balcão/administração \\
Destinatário & notificação própria & UID do caminho = destinatário & Rules/backend & lê; marcar lida só servidor \\
Sistema & trigger, job, notificação, cache & IAM mínimo; entrada server-owned & backend & sem papel humano, com dedup M7 \\
\bottomrule
\end{longtable}
\normalsize

\paragraph{Leitura e escrita.}
Catálogo de reagentes é leitura interna para papéis permitidos na Seção~3; não
expõe retirantes, auditoria ou dados de terceiros. Empréstimo completo é
legível pelo retirante, Gestor\_Almoxarifado vinculado e Chefe; bem, requisição
e histórico seguem a matriz. Consultas que poderiam retornar documento não
autorizado são negadas, pois Rules não são filtros. Escritas de domínio,
histórico, auditoria, papéis, vínculos, \texttt{Operacoes}, locks, estoques
configurados, cache, rate limit e notificações são server-owned. A única
mutação do destinatário é pedir ao backend que marque a notificação lida;
Limpar tudo não apaga documentos.

\paragraph{Claims, revogação e sessão.}
Custom Claims são projeção assinada de \texttt{roles} e
\texttt{versao\_permissoes}, gravada apenas pelo Admin SDK após o commit. O
Firebase propaga claim alterada somente para ID token novo; ID tokens têm vida
curta de cerca de uma hora. Revogar refresh token e pedir refresh melhoram a
sessão, mas não tornam claim antiga fonte autoritativa. Para leitura direta,
Rules exigem token com versão igual, Usuário ativo e relação persistida
aplicável. Para mutação crítica, backend lê documentos atuais dentro da
transação; não confia somente em claims, inclusive de Chefe\_Geral. Alteração
de papel, desativação ou remoção de vínculo incrementa versão na transação e
nega a próxima tentativa que consulte estado persistido.

\paragraph{TOCTOU, transação e idempotência.}
Para operação crítica, a transação lê Usuário, papéis, vínculo/ownership,
recurso, estado de domínio e registro M7 antes de escrever. Se revogação ou
vínculo mudar depois de uma leitura, o Firestore reexecuta a transação e a
decisão é refeita; commit posterior à revogação é negado. Verificação antes da
transação é só UX. Retry M7 concluído devolve receipt existente ao mesmo UID
sem reaplicar efeito; não permite novo efeito após perda de papel. Etapa
pendente ou efeito externo revalida antes de persistência privilegiada.

\paragraph{Security Rules, Admin SDK e UI.}
Rules protegem SDKs cliente; Admin SDK as ignora e deve aplicar esta matriz no
backend. Rules não substituem validação de operação, transação, estado de
domínio ou IAM. UI pode ocultar ação e pedir refresh, mas é UX. Collections
internas e server-owned permanecem deny-all ao cliente. Storage exige a mesma
decisão atual de papel/ownership/escopo, além de caminho, tipo, tamanho e
objeto validado.

\paragraph{Casos de regressão M9.}
Cobrir: não autenticado; sem papel; papel desconhecido; Usuário inativo;
gestor no almoxarifado correto e em outro; vínculo removido; claim antiga após
revogação; troca de papel; Professor em recurso próprio e alheio; Aluno em
turma permitida e em ação administrativa; Chefe nos limites da matriz; escrita
direta server-owned; Admin SDK sem validação; operação iniciada antes e commit
depois de revogação; retry M7 após mudança autorizativa; cache hit M8 sem
autorização; notificação de terceiro e alteração cliente-side; e tentativas de
contornar quarentena M5 ou pendência metrológica M6. Cada caso resulta em
negação ou no efeito específico da matriz, sem autorização implícita.
\subsection{Patrimônio (M10)}
\label{sec:regras-patrimonio-m10}

\begin{regra}[Autoridade, identidade e mutações patrimoniais]
\textit{Bem\_Patrimonial} é a unidade física; \textit{Resumo\_Bem\_Patrimonial}
é somente o modelo catalográfico compartilhado. O nome canônico pertence ao
resumo. \textit{nome\_equipamento} no bem e \textit{predio}/\textit{andar}/
\textit{sala} são projeções de consulta, jamais autoridades concorrentes. Uma
alteração individual de nome é reclassificação para resumo existente ou criação
de resumo novo; não renomeia o resumo de outros bens. Foto pertence apenas à
unidade física.

Toda mutação usa a ordem M9 (Auth, usuário ativo, papel/versão persistidos,
escopo global patrimonial V1 e revalidação dentro da transação) e M7
(\texttt{idOperacao}, \texttt{tipo\_operacao}, payload canônico e receipt). A
autorização é necessária, mas não dispensa unicidade, lock, versão, estado,
arquivo validado ou histórico. Gestor de Bens e Chefe podem atuar no domínio
global; Professor cria apenas requisições permitidas, nunca aprova a própria
requisição, não escreve Bem diretamente e não realiza baixa. UI e Rules são
camadas de UX/leitura; a decisão e a escrita crítica são server-owned.
\end{regra}

\begin{regra}[Plaqueta, locks e requisições]
Use \texttt{N(s)=s.trim().toUpperCase()} como única canonicalização de plaqueta.
O valor vazio ou maior que 30 caracteres após N é inválido. Bem, requisição,
lock \texttt{bem\_adicao\_\{N\}}, chave permanente
\texttt{Chaves\_Unicas/Bem\_Patrimonial\_\_\{N\}} e busca exata usam N; a
chave permanente é criada com a aprovação e jamais é liberada pela baixa.
Lock de edição é \texttt{bem\_edicao\_\{idBem\}}. Ele é criado junto da
requisição, não expira por idade e só é liberado pelo desfecho terminal que
confirma seu proprietário. Lock limita pendência; chave única protege existência
permanente; versão detecta edição concorrente; idempotência reapresenta o mesmo
resultado ao mesmo comando. Não são substitutos entre si.

Requisição de edição captura \texttt{versao\_bem\_origem}; aprovação relê o Bem
na transação e divergência termina a requisição como rejeitada por conflito,
liberando apenas seu lock. Referência de Local ou Resumo que deixou de existir é
rejeição sistêmica terminal, também com liberação vinculada. Falha transitória de
infraestrutura aborta sem transformar a requisição em rejeitada. Aprovação de
adição relê requisição, lock, chave única, Local, Resumo e foto antes de gravar,
no mesmo commit, Bem, chave, eventual resumo, resposta, histórico de cadastro e
liberação do lock.
\end{regra}

\begin{regra}[Estado, baixa, arquivo e histórico]
Conservação BOM/REGULAR/RUIM não deriva status. O ciclo V1 é
\texttt{Ativo -> Inservivel -> Ja\_dado\_baixa}; não há reversão nem salto
direto para baixa. Baixa é operação própria do Gestor/Chefe após rito SEI: exige
Bem ainda Inservivel, justificativa, PDF existente no caminho autorizado,
proprietário/origem, tamanho, assinatura \texttt{\%PDF-} e geração imutável
verificados pelo backend antes da transação. A transação grava URL/geração do
comprovante, status terminal, versão incrementada, evento \texttt{baixa} e
receipt M7. Antes da baixa, a URL é nula; depois, não pode ser reutilizada,
substituída ou apagada silenciosamente.

\texttt{Bem\_Patrimonial/\{id\}/Historico\_Patrimonio/\{eventoId\}} é a única
arquitetura normativa de histórico; coleção raiz histórica é legado divergente.
Eventos são \texttt{cadastro}, \texttt{edicao} e \texttt{baixa}; edição usa
deltas \texttt{campo, valor\_anterior, valor\_novo} do estado efetivamente
substituído. Bem e evento são atômicos; eventos usam ID determinístico da
operação/evento para retry não duplicar fatos. Snapshots de Local no evento são
imutáveis e fan-out derivado não cria histórico de negócio nem incrementa versão.
\end{regra}

\paragraph{Casos de regressão M10.}
Cobrir: criação válida, plaqueta existente e variantes de caixa/espaço, duas
adições concorrentes, resumo/local/foto ausentes ou inválidos; edição, resposta
repetida, conflito, reclassificação, local e foto; proprietário de lock errado,
rejeição/conflito liberando lock e idade sem expiração; cada transição de status;
baixa sem PDF/PDF inválido/Professor/duplicada; histórico de cadastro, delta e
baixa; retry M7, revogação M9 antes do commit, claim antiga e escrita direta de
cliente. Esses casos são critérios normativos e não atestam a implementação.

% M11 -- fonte normativa de Turmas, matrícula e convites
\subsection{Turmas, matrícula e convites (M11)}
\label{sec:regras-turmas-m11}

Esta subseção é a fonte normativa de Turma, vínculo Aluno--Turma, matrícula e
convite de ingresso (RF17, RF18, RN-TUR-01 e UI-10). Posts, comentários,
roteiros e notificações acadêmicas entram apenas nas fronteiras afetadas por
essas operações; a formalização completa desses domínios permanece fora de
M11. A autorização continua regida por M9 (Seção~\ref{sec:m9-autorizacao}), a
idempotência por M7 e a unicidade por \texttt{Chaves\_Unicas}; nada nesta
subseção cria bypass desses contratos.

\begin{regra}[Identidade e ciclo de vida da Turma]
\textit{Turma} é a entidade raiz \texttt{Turma/\{id\}}, com \texttt{id} gerado
pelo servidor. \texttt{id\_professor} é o professor dono e \texttt{id\_materia}
refere uma matéria existente, verificada no servidor, que não é recriada.
\texttt{nome\_turma} tem no máximo 100 caracteres; \texttt{ano} é inteiro;
\texttt{semestre} é exclusivamente 1 ou 2; \texttt{capacidade} é inteiro
positivo; \texttt{codigo\_turma} é único, com até 20 caracteres, gerado
exclusivamente pelo servidor; \texttt{status} pertence a \{Ativo, Arquivada\};
\texttt{data\_criacao}, \texttt{versao} (inteiro positivo) e \texttt{qtd\_alunos}
(inteiro não negativo) completam o registro.

A criação ordinária é feita pelo professor autenticado para si
(\texttt{id\_professor} = UID). Criar turma em nome de outro professor é
intervenção excepcional M9/Q13, auditada, com o dono registrado explicitamente;
a autoria não é transferida a quem não é o professor dono. Só o professor dono,
em regra, lê e altera a turma; o Chefe lê e intervém apenas nos limites de Q13.
Aluno de turma alheia não lê a turma.

Arquivar e desarquivar alteram \texttt{status} e incrementam \texttt{versao},
preservando \texttt{id}, membros, posts, comentários, histórico,
\texttt{qtd\_alunos} e \texttt{codigo\_turma}; arquivar não libera o código para
reuso. Arquivar com alunos matriculados é permitido e não é desligamento:
todos os vínculos canônicos permanecem. Turma arquivada é somente leitura
(Q08): o backend e as Rules rejeitam novo ingresso, novo convite de turma, novo
post e novo comentário, sem apagar dados. A inclusão ou remoção de membro
altera \texttt{qtd\_alunos}, vínculo e histórico, mas não é edição dos metadados
da turma e não incrementa \texttt{versao} por si; somente criação,
arquivamento/desarquivamento e edição de metadados o fazem.
\end{regra}

\begin{regra}[Matrícula canônica, espelho e contador]
\texttt{Turma/\{id\}/Alunos/\{uid\}} é o vínculo canônico e a única autoridade
de matrícula, presença e acesso acadêmico. \texttt{Usuarios/\{uid\}/Turmas/\{id\}}
é projeção de consulta que alimenta ``Minhas Turmas'' em tempo real; ela nunca
autoriza e nunca substitui o vínculo canônico (M9).

\texttt{qtd\_alunos} é o contador transacional do vínculo ativo. A leitura 3FN
\texttt{COUNT(Aluno\_x\_Turma WHERE id\_turma = X)} e \texttt{qtd\_alunos}
representam a mesma grandeza; na V1 a autoridade de vagas é
\texttt{qtd\_alunos}, lido e escrito na mesma transação que cria ou remove o
vínculo. Consultar a subcoleção \texttt{Alunos} ou recontar vínculos fora da
transação não é autoridade de vaga, pelo custo O(N) e pela concorrência.
Divergência entre \texttt{qtd\_alunos} e o número de vínculos canônicos é
inconsistência a reconciliar, nunca vaga adicional.

Ingresso, remoção e restauração de acesso atualizam, numa única transação, o
vínculo canônico, o espelho, \texttt{qtd\_alunos}, o evento em
\texttt{Historico\_Alunos\_Turma}, a auditoria e o receipt M7. Evento,
notificação ou espelho não criam matrícula por si. Nome e foto vistos por
colegas são projeções mínimas, sem matrícula nem e-mail.

Remover aluno: o professor dono desvincula (Chefe apenas por Q13). A transação
exclui o vínculo canônico e o espelho, grava \texttt{exclusao\_aluno} com
\texttt{removido\_por} e decrementa \texttt{qtd\_alunos} somente se o vínculo
existia, nunca abaixo de zero. Posts, comentários e autoria histórica são
preservados. Após remoção, o ingresso por código é bloqueado; o reingresso
exige convite explícito. Leitura de dado protegido após a remoção consulta o
vínculo canônico atual; sessão, claim ou cache antigos não restauram acesso.

Repetição (M7): reingressar quando já matriculado devolve o resultado
persistido e o acesso à turma, sem duplicar vínculo nem contagem; remover duas
vezes é idempotente e a segunda tentativa não gera decremento extra. O perdedor
de uma corrida recebe erro recuperável, nunca matrícula parcial.
\end{regra}

\begin{regra}[Capacidade, concorrência e exceção nominal]
O ingresso ordinário, por código ou por convite nominal sem exceção, exige no
commit da transação \texttt{qtd\_alunos < capacidade}. Na corrida pela última
vaga, o Firestore reexecuta a transação, um único commit cria o vínculo e o
perdedor recebe \texttt{failed-precondition} sem matrícula parcial.

Somente convite nominal válido, emitido pelo professor dono, com
\texttt{exceder\_capacidade = true} e \texttt{justificativa\_excecao}
persistida, pode ultrapassar a capacidade; a exceção é revalidada na aceitação.
Ela é específica daquele convite e daquela identidade, não é estado global da
turma nem passe livre por posse do código. Convite sem exceção continua sujeito
a vaga, e o código de turma nunca autoriza exceção.

Redução de capacidade: é proibido reduzir \texttt{capacidade} para valor menor
que \texttt{qtd\_alunos}; o servidor rejeita, fail-closed, com erro recuperável
e orienta remover membros antes ou manter a capacidade. Assim não há expulsão
implícita de alunos nem vaga negativa, e a capacidade nunca fica abaixo da
ocupação. \texttt{qtd\_alunos} nunca é negativo e pode ser reconciliado com os
vínculos canônicos.
\end{regra}

\begin{regra}[Convites: token, unicidade e fronteira externa]
\texttt{Convite\_Aluno/\{id\}} é coleção raiz; \texttt{id\_turma} \texttt{NULL}
identifica convite global. O e-mail é armazenado normalizado; o token aleatório
de 32 bytes (CSPRNG) é guardado apenas como hash SHA-256, comparado em tempo
constante, com expiração de sete dias, status \texttt{pendente}, \texttt{aceitado}
ou \texttt{expirado}, \texttt{convidado\_por}, \texttt{convidado\_em},
\texttt{numero\_matricula} opcional, \texttt{exceder\_capacidade} e
\texttt{justificativa\_excecao}, \texttt{aceitado\_por} e \texttt{aceitado\_em}.
O docId é determinístico, derivado no servidor por HMAC-SHA-256 do contexto
(turma ou global) e do e-mail normalizado com segredo do servidor; nunca se usa
hash simples e previsível de e-mail como chave nem se expõe o e-mail em claro.

Unicidade: no máximo um convite \texttt{pendente} por
(\texttt{email} normalizado, \texttt{id\_turma}); para convite global, no
máximo um pendente por e-mail normalizado com \texttt{id\_turma} \texttt{NULL}.
A chave determinística em \texttt{Chaves\_Unicas} distingue as duas situações e
é validada na transação. Convite aceito ou expirado não ocupa a pendência.
Reenvio substitui hash e expiração no mesmo documento e invalida o link
anterior. Convite de turma e convite global são distintos: aceitar convite de
turma cria vínculo, espelho e contador; aceitar convite global não cria
matrícula nem contador, apenas garante \texttt{Usuarios} e o papel Aluno, e o
aluno ingressa depois por código ou por outro convite. Matrícula é obrigatória
e única na identidade final \texttt{Aluno}, opcional e não única no convite.

Expiração operacional: \texttt{expira\_em} vencido torna o convite inaceitável
e ele deixa de contar como pendente para a unicidade; job ou reconciliação
idempotente materializa \texttt{expirado}. Toda pendência termina por aceitação,
expiração ou reenvio, nunca fica eterna.

Fronteira Auth, Firestore e e-mail: a aceitação exige sessão com e-mail
verificado correspondente. Criação condicional de \texttt{Usuario},
\texttt{Aluno}, papel e vínculo e o consumo do convite ocorrem na mesma
transação Firestore; a criação de conta Firebase Auth quando inexistente e o
envio de e-mail são etapas externas, pós-commit, recuperáveis e idempotentes.
Falha parcial não pode criar conta, papel, vínculo ou espelho divergente nem
consumir o convite sem os efeitos correspondentes. Criar o registro do convite
não é enviar o e-mail.
\end{regra}

\begin{regra}[Aceitação, reingresso e composição M7/M9]
Aceitar convite valida, na mesma transação: convite pendente não expirado,
e-mail verificado, identidade, status Ativo da turma, vaga ou exceção válida e
ausência de vínculo ativo. A aceitação é única e idempotente (M7): repetir a
aceitação devolve o resultado persistido, sem recriar vínculo nem incrementar
\texttt{qtd\_alunos}. Convite consumido, expirado, de e-mail divergente ou de
turma arquivada resulta em negação específica e recuperável. Aluno removido
anteriormente só reingressa por convite explícito, nunca pelo código.

Toda mutação M11 lê autorização e recurso na mesma transação (M9): professor
dono, conta ativa e versão de permissões, com TOCTOU coberto pela reexecução da
transação; operação iniciada antes de revogação não grava depois dela. UI-10 e
a Seção~\ref{sec:fluxos-completos} são UX: botões, código e navegação não
autorizam. O frontend não afirma envio de e-mail quando apenas o registro foi
criado e distingue carregamento, sucesso, vazio, erro e negação.
\end{regra}

\paragraph{Casos de regressão M11.}
Cobrir: corrida pela última vaga; dois convites simultâneos para o mesmo
(e-mail, turma) e para o global; replay de aceite; remoção concorrente e dupla;
arquivamento com ocupação, desarquivamento e escrita em turma arquivada; convite
global sem matrícula; ingresso de removido por código (negado) e por convite
(exceção válida); exceção de capacidade com e sem justificativa; redução de
capacidade abaixo da ocupação; contador nunca negativo; espelho ou cache
desatualizado não autoriza; retry M7 e revogação M9 antes do commit. Esses
casos são critérios normativos e não atestam a implementação.

\subsubsection{Seleção do tipo em Resumo de Reagente}
\begin{regra}
Ao cadastrar uma \textit{Especificacao\_Reagente}, o estado físico do resumo (\texttt{SOLIDO} ou \texttt{LIQUIDO}) determina quais campos operacionais de frasco serão apresentados. O tipo de substância (\texttt{PURA} ou \texttt{MISTURA}) determina as propriedades químicas obrigatórias, conforme a matriz de obrigatoriedade. DP-A01: estado físico e higroscopicidade pertencem ao resumo; unidade é derivada e densidade permanece na especificação. \\
\end{regra}
\subsubsection{Excluindo um Aluno de uma turma}
\begin{regra}
Quando o professor exclui um aluno de uma turma, ele não apga o aluno d sistema. Apenas o retira da turma. Seu histórico nos posts é mantido. Ou seja, o alunos que ainda estão na turma conseguem ver os comentários dele e também seu nome e foto de perfil nos comentários dele.
\end{regra}
\subsubsection{Chefe Geral}
\begin{regra}
O sistema pode ter mais de um chefe geral, isso não altera o sistema.
\end{regra}
\subsubsection{O responsável pelo bem patrimonial não é um Usuário do sistema LCQUI}
\begin{regra}
Por iss que "\textit{nome\_responsavel\_sei}" é um simples texto.
\end{regra}
\subsubsection{Historico de Posts na Turma}
\begin{regra}
A entidade \textit{Historico\_Posts\_Turma} registra no campo \textit{editado\_por} o identificador do usuário que realizou a edição. Normalmente é o Professor dono da turma, mas o Chefe Geral também pode editar ou remover posts em caráter excepcional de moderação institucional (Q13); nesses casos o campo \textit{editado\_por} identifica a ação da chefia e o campo \textit{acao} em \textit{Registro\_de\_Auditoria} registra a motivação. O \textit{id\_turma} não é armazenado nesta entidade 3FN porque pode ser obtido por \textit{Historico\_Posts\_Turma.id\_post} $\rightarrow$ \textit{Post.id\_turma}; se a consulta direta por turma for necessária no Firestore, esse vínculo poderá ser denormalizado na coleção de histórico.
\end{regra}
\subsubsection{Não compartilhar o mesmo roteiro 2 vezes com um professor}
\begin{regra}
Na tabela: \textit{Roteiro Professor Compartilhado} \\
UNIQUE(id\_roteiro\_experimento, id\_professor)
\end{regra}
\subsubsection{Densidade e concentração na especificação do reagente}
\begin{regra}
A densidade é uma propriedade da \textit{Especificacao\_Reagente}, e não do \textit{Resumo\_Reagente}. Ela deve ser informada quando obrigatória para a especificação e, após a confirmação cadastral, não pode ser alterada diretamente, pois participa dos cálculos de volume e consumo dos frascos. A concentração pertence à composição da especificação. Duas especificações com o mesmo nome comercial podem ter densidade, concentração, fabricante ou código de produto diferentes e, portanto, são especificações distintas sob o mesmo resumo quando ainda representam o mesmo item químico agrupador; não se deve criar um novo \textit{Resumo\_Reagente} apenas por haver diferença de concentração ou densidade.
\end{regra}
\subsubsection{Regra para código do frasco de Reagente}
\begin{regra}
O código do frasco de reagente (formato \texttt{LCQUI-N}) é gerado e garantido exclusivamente no momento do cadastro do frasco via transação atômica. A impressão de etiquetas virgens prévias fornece identificadores sugeridos para colar fisicamente, mas não cria frascos no sistema nem garante a reserva do código.
\end{regra}

\subsubsection{Segurança na Reimpressão de Frascos Cadastrados}
\begin{regra}
Para evitar a clonagem acidental de etiquetas na bancada e risco de trocas químicas, a reimpressão de segunda via é limitada a no máximo 10 frascos por sessão. O documento impresso deve conter a Ficha de Conferência e Rastreabilidade completa e 1 única etiqueta de reposição por frasco.
\end{regra}
\subsubsection{Regra em Notificação para Professor}
\begin{regra}
O campo \textit{id\_turma} é obrigatório para todas as notificações de contexto acadêmico (\texttt{COMENTARIO}, \texttt{POST}, \texttt{ADICIONADO}, \texttt{REMOVIDO}, \texttt{TURMA\_ARQUIVADA}, \texttt{TURMA\_DESARQUIVADA}) e deve ser mantido estritamente \texttt{NULL} para notificações operacionais e de almoxarifado/patrimônio. A notificação \texttt{DATA\_DEVOLUCAO\_REAGENTE} é preventiva: deve ser emitida pelo backend quando um empréstimo ainda está \texttt{EM\_USO} e sua data de devolução é hoje ou amanhã. A detecção é responsabilidade da Cloud Function; o frontend apenas apresenta e pode agrupar visualmente os registros em ``atrasados'', ``vence hoje'' e ``vence amanhã''. O frontend não decide se uma notificação deve ser criada.

O job pode usar uma única consulta por \texttt{status = EM\_USO} e \texttt{data\_devolucao\_prevista <= fim\_de\_amanha}, mas deve classificar cada empréstimo no servidor antes de criar a notificação, separando a janela de atraso da janela preventiva e aplicando idempotência. A idempotência é garantida por \textbf{docId determinístico} composto por \texttt{\{id\_emprestimo\}--\{janela\}}, onde \texttt{janela} é \texttt{VENCE\_HOJE} ou \texttt{VENCE\_AMANHA}; a Cloud Function usa \texttt{set(..., \{merge: false\})} apenas se o documento ainda não existir, evitando duplicatas em execuções repetidas do job diário. Um empréstimo que muda de janela (amanhã $\rightarrow$ hoje) cria um segundo documento; o anterior não é removido, apenas expirado.

Deve aparecer na interface:
\begin{itemize}
    \item COMENTÁRIO: x novos comentários de Fulano de Tal na turma y;
    \item REQUISICAO\_BEM: resposta de aprovação ou rejeição enviada ao Professor solicitante;
    \item REQUISICAO\_ADICAO\_BEM: nova solicitação de adição enviada aos Gestores de Bens Patrimoniais;
    \item REQUISICAO\_EDICAO\_BEM: nova solicitação de edição enviada aos Gestores de Bens Patrimoniais;
    \item DATA\_DEVOLUCAO\_REAGENTE: a data de devolução de um ou mais frasco(s) está expirando [amanhã, hoje];
    \item ROTEIRO\_COMPARTILHADO: Você tem x novo(s) roteiro(s) compartilhado(s) com você.
\end{itemize}
\end{regra}

\subsubsection{Regra para convidar um aluno que não está no sistema}
\begin{regra}
Quando o aluno é convidado sem especificar uma turma, id\_turma é NULL em \textit{Convite para Aluno}. Quando o professor vai convidar o aluno, ele pode ou não especificar uma turma, mas ao enviar um convite, tem que aparecer uma confirmação no caso de estar convidando um aluno sem uma turma para ele.
\end{regra}
\subsubsection{Regra para registrar um patrimônio como dado baixa}
\begin{regra}
É preciso fazer upload do documento pdf que comprove que o bem patrimonial foi dado baixa no sistema da UENF
\end{regra}
\subsubsection{Cálculo de Volume e Peso de Reagentes (Padrão de Laboratório)}
\begin{regra}[Controle de Volume via Pesagem]
O sistema adotará a pesagem como método principal de controle de volume dos frascos, utilizando a fórmula $Volume = \frac{Peso_{total} - Peso_{vazio}}{Densidade}$, em que $Peso_{vazio}$ é a tara real do recipiente, obtida pela medição do recipiente efetivamente vazio.

\textbf{1. Cadastro de Frasco Fechado:}\\
O usuário informa o Peso Total (frasco + reagente) medido na balança e o Conteúdo Nominal do rótulo (volume em mL para líquidos, massa em g para sólidos). O sistema deriva a tara teórica do recipiente a partir do peso total e do conteúdo nominal, quando disponível. Essa tara é apenas uma referência derivada enquanto o frasco contém produto e não é medição física definitiva do recipiente vazio:\\
\begin{itemize}
    \item Reagente Líquido ($mL$): $Peso_{vazio} = Peso_{total} - (Volume_{nominal} \times Densidade)$
    \item Reagente Sólido ($g$): $Peso_{vazio} = Peso_{total} - Massa_{nominal}$
\end{itemize}
O frasco inicia com \texttt{condicao\_inicial\_cadastro = FECHADO}, \texttt{estado\_fisico\_frasco = FECHADO} e saldo conhecido somente com dados quantitativos suficientes. Um cadastro FECHADO admite conteúdo nominal NULL quando desconhecido ou ilegível, nunca zero. Nesse caso, peso\_frasco\_vazio e origem\_tara são NULL e saldo\_desconhecido é true. Abertura ou pesagem bruta isolada não resolvem esse desconhecimento. Sem tara real, segue o mesmo ciclo de vida de JA\_ABERTO em situação metrológica equivalente: não exigir tara fabricada para quebra, descarte, extravio ou reencontro. EXTRAVIADO preserva a flag; \texttt{VAZIO} e \texttt{QUEBRADO} são estados fisicamente não utilizáveis e indisponíveis, normalmente encaminhados para descarte, mas \textbf{não} são estados terminais da máquina (\texttt{DESCARTADO} é o estado terminal); nada disso inventa peso ou quantidade química.

\textbf{2. Cadastro de Frasco Já Aberto:}\\
O frasco inicia com \texttt{condicao\_inicial\_cadastro = JA\_ABERTO}, \texttt{estado\_fisico\_frasco = ABERTO}, \texttt{peso\_frasco\_vazio = NULL} e \texttt{saldo\_desconhecido = true}. Não existem as modalidades de tara conhecida, estimativa de volume ou estimativa de massa, nem campos de peso do frasco vazio informado, volume atual estimado ou massa atual estimada. O usuário informa apenas dados reais disponíveis: o peso total atual medido na balança e, quando conhecido ou legível, o conteúdo nominal do rótulo. O conteúdo nominal é declaração original do fabricante e não representa o saldo atual; quando desconhecido ou ilegível, permanece \texttt{NULL}, nunca zero. Enquanto o recipiente não for esvaziado e pesado vazio, o saldo atual permanece desconhecido e a tara do recipiente ainda não é conhecida; não se estima saldo inicial.

\textbf{3. Empréstimo e Devolução:}\\
No empréstimo, registra-se o \textit{peso\_saida}. Na devolução, registra-se o \textit{peso\_retorno}. Quando quantitativamente validado, o consumo é documentado na tabela \textit{Emprestimo\_Reagente} como: $Volume_{utilizado} = \frac{\max(0,\ (Peso_{saida} - Peso_{retorno}^{efetivo}) - Perda_{evaporacao})}{Densidade}$ (líquidos); sólidos usam a diferença não negativa em g. A perda por evaporação é contabilizada separadamente do consumo (campo e evento \texttt{EVAPORACAO}); ganho tolerado segue Q06, com consumo zero e ajuste separado.

\textbf{4. Esgotamento e Tara Real:}\\
A determinação de que o recipiente ficou vazio é confirmada explicitamente pelo gestor na devolução, tanto para frascos originalmente \texttt{FECHADO} quanto para frascos \texttt{JA\_ABERTO}, por meio da decisão ``O frasco ficou vazio nesta devolução? [ ] Sim''. A comparação com a tara anterior não é autoridade exclusiva para declarar esgotamento.

Sem tara real anterior, ao confirmar o esgotamento, o peso aferido na devolução (\textit{peso\_retorno}) passa a ser a \textbf{tara real} do recipiente: \texttt{peso\_frasco\_vazio = peso\_retorno}, \texttt{peso\_atual = peso\_retorno}, \texttt{estado\_fisico\_frasco = VAZIO}, \texttt{disponibilidade = INDISPONIVEL} e \texttt{saldo\_desconhecido = false}. Não se altera o \texttt{conteudo\_nominal} e registra-se o evento \texttt{FICOU\_VAZIO}. Para frasco \texttt{JA\_ABERTO}, \textit{peso\_retorno} é a primeira tara real conhecida. Não existe recalibração parcial por estimativa: recalibrar tara significa recipiente efetivamente vazio com peso vazio real medido e justificativa humana conforme a lista fechada; não existe correção administrativa metrológica na V1.

\textbf{Tolerância de Devolução (Regra Q06 - PDF-025):}\\
A margem de tolerância é estritamente baseada no \textbf{peso bruto de saída} ($Peso_{saida}$). Rejeita-se explicitamente o cálculo sobre massa líquida, pois a variância incide sobre a totalidade do corpo medido na balança (incluindo o frasco).
\begin{itemize}
    \item Reagente normal: $\max(1\,g, 0{,}005 \times Peso_{saida})$
    \item Reagente higroscópico: $\max(2\,g, 0{,}02 \times Peso_{saida})$
\end{itemize}
A avaliação de ganho máximo permitido segue estritamente a fórmula: $Peso_{retorno} > Peso_{saida} + Tolerancia$. 
Um ganho de massa acima dessa tolerância \textbf{não aborta a operação via rollback}. O sistema conclui a devolução em modalidade extraordinária: o empréstimo transiciona para \texttt{DEVOLVIDO\_COM\_ANOMALIA}, e o frasco transiciona compulsoriamente para \texttt{disponibilidade = INDISPONIVEL} e \texttt{em\_quarentena = true}, com \texttt{detalhe\_status} registrando suspeita de contaminação/adulteração. É terminantemente proibido registrar evento de higroscopia (\texttt{ganho\_massa\_higroscopia}) para frascos não-higroscópicos ou com ganhos fora da tolerância legal.

Um retorno abaixo da tara real conhecida é preservado exatamente como medido; não existe limiar fixo adicional. Sem confirmação de vazio, encerrar a custódia como DEVOLVIDO\_COM\_ANOMALIA, bloquear o frasco em quarentena/INDISPONIVEL e manter consumo final pendente. A referência teórica não é limite físico absoluto. Aplicam-se os contratos de resolução da Seção 10; Q06 compara retorno com saída, não retorno com tara.
\end{regra}

\subsubsection{Exclusividade na retirada de reagentes}
\begin{regra}
O Chefe Geral atua como operador de balcão, nunca como destinatário/retirante. Apenas os usuários cadastrados como \textit{Professor} ou \textit{Bolsista} têm permissão para retirar reagentes do almoxarifado (o \textit{Emprestimo\_Reagente.id\_usuario\_retirou} deve corresponder a um \textit{Usuario} com uma dessas duas linhas associadas). Demais Alunos (mesmo de iniciação científica ou TCC, sem o papel de Bolsista) não possuem autorização no sistema para esta ação. Esta regra foi corrigida para não contradizer a seção \textit{Bolsista} (3.6), que descreve explicitamente o gestor de reagentes fazendo empréstimos para usuários Bolsista. O autoatendimento é exclusivo de usuários portadores da combinação \texttt{Professor + Gestor\_Almoxarifado}. Usuários no papel \texttt{Aluno + Gestor\_Almoxarifado} são terminantemente impedidos pelo backend de retirar reagentes para si mesmos, em respeito à matriz de segregação de funções (SOD).
\end{regra}

\subsubsection{Usuário Multi-role, exceto Chefe Geral}
\begin{regra}
Uma conta que possui o papel \textit{Chefe Geral} não pode possuir nenhum outro papel. Se a mesma pessoa exercer, na instituição, a função de Chefe Geral e também uma função representada por outro papel do sistema, deverão existir duas contas independentes, cada uma com identidade e permissões próprias.
\end{regra}
\subsubsection{Usuário Multi-role, Aluno não pode ser professor}
\begin{regra}
O usuário que é aluno não pode ser professor.
\end{regra}
\subsubsection{Usuário Multi-role, Quais multi papeis podem existir}
\begin{regra}
Pode haver:\\
Professor que é Gestor de Bens Patrimoniais.\\
Professor que é Gestor de Almoxarifado\\
Professor que é Gestor de Bens Patrimoniais e Gestor de Almoxarifado\\
Aluno que é Bolsista\\
Aluno que é Bolsista e Gestor de Bens Patrimoniais\\
Aluno que é Gestor de Almoxarifado\\
Aluno que é Gestor de Bens Patrimoniais\\
Aluno que é Gestor de Bens Patrimoniais e Gestor de Almoxarifado\\
O Usuário Aluno que também é Bolsista não pode ser Gestor de Almoxarifado. Se, na vida real, um aluno é bolsista e também precisa ser Gestor de Almoxarifado, devem ser criadas duas contas para ele no sistema.\\
No caso do Aluno que exerce algum papel de Gestor (ou ambos): o aluno pode editar bens patrimoniais e/ou reagentes, de acordo com seus papéis respectivos, porque também tem outro papel além de Aluno.
\end{regra}
\subsubsection{Aluno que é bolsista}
\begin{regra}
Apenas o chefe geral consegue colocar um aluno como bolsista.
\end{regra}
\subsubsection{Status, vencimento e descarte de frascos}
\begin{regra}
\begin{itemize}
    \item Um frasco \texttt{VAZIO} ou \texttt{QUEBRADO} deve ser exibido como ``Pendente de Descarte'', salvo se já estiver \texttt{DESCARTADO}.
    \item Um frasco \texttt{VENCIDO} não é automaticamente descartado e também não é automaticamente bloqueado para uso. Para frasco vencido, o sistema mantém a decisão explícita: \texttt{em\_quarentena = true}, ou \texttt{uso\_vencido\_autorizado = true}, ou situação de \textit{PENDENTE DE DESCARTE} com \texttt{uso\_vencido\_autorizado = false}.
    \item Em cadastro, quando a validade efetiva já tiver expirado no instante da operação, o gestor é obrigado a escolher um desses três destinos: \texttt{QUARENTENA}, \texttt{PENDENTE\_DE\_DESCARTE} ou \texttt{DISPONIVEL} com uso vencido explicitamente autorizado conforme Q04.
    \item O frasco em quarentena não pode ser retirado; o frasco pendente de descarte não pode ser retirado; o frasco vencido com uso explicitamente autorizado pode ser retirado, mas a interface deve alertar transparentemente antes da confirmação do empréstimo.
    \item Se um frasco vencer enquanto estiver emprestado, o empréstimo continua válido. A devolução \textbf{não} recalcula o vencimento com a hora atual: a passagem do tempo é responsabilidade do job de vencimento, e a devolução lê o estado canônico persistido \textit{Frasco\_Reagente.vencido}. O snapshot imutável \texttt{vencido\_na\_retirada} distingue ``venceu durante o empréstimo'' (não estava vencido na retirada) de ``já estava vencido na retirada''; validade desconhecida sem vencimento afirmativo é o terceiro caso. Em qualquer desses casos o gestor registra o destino pós-validade.
    \item A ação de descarte institucional (\texttt{descartarFrasco}) exige \texttt{em\_quarentena = false} e \texttt{disponibilidade != EMPRESTADO}, e é permitida para: (1) frascos com estado físico \texttt{VAZIO}, (2) frascos \texttt{QUEBRADO}, (3) frascos com \texttt{vencido = true} e \texttt{uso\_vencido\_autorizado = false}, e (4) frascos com autorização técnica de descarte ativa. A autorização técnica da rota (4) é produzida exclusivamente pela decisão humana \texttt{PENDENTE\_DE\_DESCARTE} de saída de quarentena e nunca por texto livre em \texttt{detalhe\_status}. \texttt{INDISPONIVEL} isolado não habilita descarte. As rotas (1)--(3) admitem conteúdo químico remanescente nos estados \texttt{ABERTO} ou \texttt{FECHADO}, formalizando a baixa do passivo químico perigoso. A quarentena deve ser resolvida antes; se \texttt{em\_quarentena = true}, nenhuma rota se aplica diretamente.
\end{itemize}
\end{regra}

\subsubsection{Extravio, reencontro e quarentena (M5)}
\begin{regra}
Consolidação normativa do ciclo de extravio, reencontro e quarentena. As operações são server-owned, exigem papel \textit{Chefe\_Geral} ou \textit{Gestor\_Almoxarifado} com vínculo ao almoxarifado do frasco e justificativa humana (mínimo de 20 caracteres após \texttt{trim}) para extravio e reencontro; a autorização completa de papéis pertence a M9.
\begin{itemize}
    \item \textbf{Extravio.} Um frasco pode ser declarado extraviado a partir de qualquer localização LOCALIZADO e estado físico, exceto \texttt{DESCARTADO} (terminal). O efeito é \texttt{situacao\_localizacao = EXTRAVIADO}, o último \texttt{estado\_fisico\_frasco} é preservado e a \texttt{disponibilidade = INDISPONIVEL}. A autorização corrente de descarte é revogada no extravio; a decisão anterior permanece no histórico auditável e não autoriza descarte automático após o reencontro. Extravio \textbf{não} é consumo, esgotamento, quebra nem peso zero: não fabrica tara, saldo, peso, consumo ou evidência metrológica. São preservados \textit{saldo}, \textit{saldo\_desconhecido}, \textit{peso\_atual}, \textit{peso\_no\_cadastrado}, tara e \textit{origem\_tara}, \textit{validade\_fechado}/\textit{validade\_efetiva}/\textit{validade\_desconhecida}, \textit{vencido}, \textit{uso\_vencido\_autorizado}, \textit{abertura\_historica\_desconhecida} e a localização. O \textit{estado físico anterior} é o último estabelecido na trilha histórica e não é apagado.
    \item \textbf{Extravio durante empréstimo ativo.} Se existir empréstimo \texttt{EM\_USO} ou \texttt{ATRASADO} para o frasco, ele é encerrado como \texttt{ENCERRADO\_EXTRAORDINARIO} com \texttt{tipo\_encerramento\_excepcional = EXTRAVIO\_SINISTRO}, \texttt{motivo\_encerramento\_excepcional} e \texttt{id\_gestor\_encerramento}. Como \textbf{não houve retorno físico}: \texttt{peso\_retorno}, \texttt{peso\_retorno\_efetivo}, \texttt{medida\_utilizada} e \texttt{id\_resolucao\_metrologica} permanecem \texttt{NULL}, \texttt{consumo\_validado = FALSE} e \texttt{data\_devolucao\_efetuada} permanece \texttt{NULL} (o instante do encerramento fica no histórico/operação; não é devolução). \texttt{peso\_saida} é preservado. \texttt{massa\_perda\_estimada\_g} só pode ser preenchida quando há tara conhecida e é uma \textbf{estimativa} gravimétrica, nunca medição nem consumo validado (fronteira M6). Empréstimo ativo que já possua \texttt{peso\_retorno} é inconsistência e exige reconciliação manual.
    \item \textbf{Reencontro.} Só um frasco cuja \texttt{situacao\_localizacao = EXTRAVIADO} pode ser reencontrado. O reencontro grava LOCALIZADO, impõe quarentena, revoga qualquer autorização corrente remanescente e preserva o último estado físico conhecido. É um \textbf{novo fato físico}: não desfaz o extravio e não reabre empréstimo encerrado. O gestor informa o estado físico constatado (\texttt{ABERTO}, \texttt{FECHADO}, \texttt{VAZIO} ou \texttt{QUEBRADO}) e o frasco passa a \texttt{em\_quarentena = TRUE} e \texttt{disponibilidade = INDISPONIVEL}. O estado constatado é validado contra o \textbf{último estado físico conhecido antes do extravio} (\texttt{estadoAntesDoExtravio}), obtido da trilha histórica ordenada sem criar coluna nova: \texttt{FECHADO} só é aceito se antes era \texttt{FECHADO} e não havia abertura conhecida; \texttt{ABERTO} só é aceito se antes era \texttt{ABERTO} ou \texttt{FECHADO} (neste caso a abertura ocorreu em período sem rastreabilidade temporal precisa); \texttt{VAZIO} só é aceito se antes já era \texttt{VAZIO} (restauração sem nova medição); \texttt{QUEBRADO} representa quebra física constatada. Transições fisicamente impossíveis são rejeitadas, como \texttt{ABERTO} $\rightarrow$ \texttt{EXTRAVIADO} $\rightarrow$ \texttt{FECHADO}, \texttt{VAZIO} $\rightarrow$ \texttt{EXTRAVIADO} $\rightarrow$ \texttt{FECHADO}/\texttt{ABERTO} e \texttt{QUEBRADO} $\rightarrow$ \texttt{EXTRAVIADO} $\rightarrow$ \texttt{ABERTO}/\texttt{FECHADO}. \texttt{saldo}, \texttt{saldo\_desconhecido}, \texttt{peso\_atual}, tara e \texttt{vencido} permanecem; uma medição válida posterior é a única forma de alterar o saldo. Se o frasco foi encontrado \texttt{ABERTO} vindo de \texttt{FECHADO} (ou sem data de abertura determinável), marca-se \texttt{abertura\_historica\_desconhecida = TRUE}, \texttt{data\_abertura = NULL}, \texttt{validade\_desconhecida = TRUE} e \texttt{validade\_efetiva = NULL}; nunca se inventa a data. Se o frasco não era previamente \texttt{VAZIO} mas parece vazio ao ser reencontrado, M5 \textbf{não} inventa esgotamento, tara ou quantidade zero: registra a observação, mantém a quarentena e a confirmação quantitativa/metrológica pertence a M6.
    \item \textbf{Validade no reencontro.} O reencontro \textbf{não} recalcula vencimento pelo relógio. A passagem do tempo é responsabilidade do job de vencimento (Seção~10.7) e o estado persistido \textit{Frasco\_Reagente.vencido} é a autoridade operacional. O reencontro apenas lê esse estado e, no caso de indeterminação exata, marca \texttt{validade\_desconhecida}; não agenda, não revalida e não antecipa as correções de validade da V2.
    \item \textbf{Quarentena.} A quarentena é uma dimensão operacional própria, distinta de \texttt{INDISPONIVEL}: um frasco em quarentena permanece \texttt{INDISPONIVEL} e nunca é liberado para nova retirada enquanto a quarentena não for resolvida. Ela é ortogonal a \textit{estado físico}, \textit{validade}, \texttt{saldo\_desconhecido} e anomalia metrológica. \texttt{vencido = TRUE} nunca é usado como mecanismo de quarentena. Enquanto durar, ficam proibidas a retirada, a devolução (não haveria empréstimo ativo) e o descarte direto.
    \item \textbf{Saídas da quarentena.} O conjunto é fechado: (i) \texttt{VOLTAR\_A\_DISPONIVEL}, que encerra a quarentena, devolve \texttt{DISPONIVEL} e exige estado físico \texttt{ABERTO}/\texttt{FECHADO}, ausência de pendência metrológica e que o frasco não esteja vencido sem uso autorizado; (ii) \texttt{PENDENTE\_DE\_DESCARTE}, que encerra a quarentena, mantém \texttt{INDISPONIVEL} e cria a autorização operacional estruturada consumida por \texttt{descartarFrasco} (HQ-M2-008/B); (iii) \textbf{permanência} em quarentena, que não é uma decisão de saída e mantém o bloqueio. \texttt{VOLTAR\_A\_DISPONIVEL} \textbf{nunca} pode liberar \texttt{VAZIO}, \texttt{QUEBRADO} ou \texttt{DESCARTADO}; para frasco reencontrado \texttt{VAZIO} ou \texttt{QUEBRADO}, a resolução só pode manter a quarentena ou produzir \texttt{PENDENTE\_DE\_DESCARTE}. Nenhuma saída revalida, renova, estende ou zera o vencimento, e nenhuma delas volta o frasco a \texttt{DISPONIVEL} por conveniência. Não existe \texttt{QUARENTENA} $\rightarrow$ \texttt{DESCARTADO} direto; o descarte técnico exige a etapa \texttt{PENDENTE\_DE\_DESCARTE}.
    \item \textbf{Rastreabilidade.} Extravio, reencontro, entrada e resolução de quarentena geram eventos imutáveis no \textit{Historico\_Frasco\_Reagente} (\texttt{EXTRAVIO}, \texttt{REENCONTRO}, \texttt{ENTROU\_EM\_QUARENTENA}, \texttt{LIBERADO\_QUARENTENA}, \texttt{PENDENTE\_DE\_DESCARTE}) com frasco, almoxarifado, ator, instante, motivo e vínculo ao empréstimo quando aplicável. O reencontro não apaga o extravio; a resolução não apaga o reencontro; nada é sobrescrito. O estado anterior é o último estabelecido na própria trilha ordenada.
\end{itemize}
As operações críticas de M5 seguem o contrato global de idempotência (Seção~\ref{sec:idempotencia-m7}): \texttt{idOperacao} obrigatório com identidade \texttt{(uid, tipo\_operacao, payload\_hash)}; retry não reaplica extravio, reencontro nem resolução de quarentena.
\end{regra}

\subsubsection{Q06, tara, medição e resolução metrológica (M6)}
\begin{regra}[Vocabulário metrológico e autoridades]
Cada campo quantitativo possui uma única autoridade. Não usar dois campos como autoridade do mesmo fato sem explicitar a diferença.
\begin{itemize}
    \item \texttt{peso\_saida}: leitura bruta da balança na retirada; imutável; base da tolerância Q06.
    \item \texttt{peso\_retorno}: observação física original da devolução; persistida e nunca reescrita por resolução posterior.
    \item \texttt{peso\_retorno\_efetivo}: valor quantitativo validado que a resolução autoriza usar no cálculo definitivo; \texttt{NULL} enquanto a pendência estiver aberta; não apaga nem substitui \texttt{peso\_retorno}.
    \item \texttt{peso\_atual}: observação corrente do frasco; não prova o peso no instante histórico da devolução.
    \item \texttt{peso\_frasco\_vazio} + \texttt{origem\_tara}: tara (\texttt{NULL}, \texttt{REFERENCIA\_TEORICA} ou \texttt{MEDIDA\_REAL}).
    \item \texttt{peso\_no\_cadastrado}: peso total inicial; imutável.
    \item \texttt{medida\_utilizada}: consumo quantitativamente validado; \texttt{NULL} \textbf{não} é zero.
    \item \texttt{peso\_perda\_evaporacao}: perda em g registrada separadamente do consumo.
    \item \texttt{massa\_perda\_estimada\_g}: estimativa de sinistro/extravio (fronteira M5); não é medição nem consumo.
    \item \texttt{saldo} (\texttt{saldo\_aferido\_g}/\texttt{saldo\_aferido\_ml}) e \texttt{saldo\_desconhecido}: saldo corrente server-owned; desconhecido \textbf{não} é zero.
    \item \texttt{consumo\_validado}: autoridade que separa resultado quantitativo encerrado de pendência aberta.
    \item \texttt{anomalia\_metrologica}: classificação textual server-owned da inconsistência.
    \item \texttt{id\_resolucao\_metrologica}: vínculo ao evento auditável que encerrou a pendência.
    \item \texttt{densidade\_aplicada}: snapshot histórico da densidade no empréstimo.
\end{itemize}
Regra geral: desconhecido $\neq$ zero, ausente $\neq$ zero, não mensurável $\neq$ zero, estimado $\neq$ medido, atual $\neq$ histórico, referência teórica $\neq$ tara real. Nunca preencher um valor apenas para completar o modelo.
\end{regra}

\begin{regra}[Q06 --- tolerância de ganho de massa]
Q06 compara \textbf{retorno com saída}, sobre o peso bruto \texttt{peso\_saida}: reagente normal $tolerancia = \max(1\,g,\ 0{,}005 \times peso_{saida})$; higroscópico $tolerancia = \max(2\,g,\ 0{,}02 \times peso_{saida})$. A condição de ganho anômalo é $peso_{retorno} > peso_{saida} + tolerancia$. O antigo limiar fixo de 5\,g está definitivamente eliminado e não pode ser recriado como definição de vazio, segunda tolerância, fallback, exceção ou justificativa. Q06 \textbf{não} é regra de tara: nunca comparar \texttt{peso\_retorno} com \texttt{peso\_frasco\_vazio}; não existe ``tolerância de tara Q06'', margem percentual para retorno abaixo da tara nem limiar fixo de vazio.
\end{regra}

\begin{regra}[Consumo, evaporação e densidade]
A perda bruta é $\Delta_{bruto} = \max(0,\ peso_{saida} - peso_{retorno}^{efetivo})$. A evaporação é registrada separadamente do consumo (campo e evento \texttt{EVAPORACAO}) e obedece à constraint $0 \le \text{Evap} \le \Delta_{bruto}$; não é autoridade de vazio nem de Q06. Se a evaporação informada exceder a perda bruta, a operação é \textbf{rejeitada} antes da conclusão quantitativa: não há clamp silencioso, conversão automática de $8\,g$ em $5\,g$, aceitação com consumo zero nem anomalia que “salve” o dado. Quando a constraint é satisfeita, $massa_{consumida} = \Delta_{bruto} - \text{Evap} \ge 0$ naturalmente; o $\max(0,\cdot)$ externo não é mecanismo para esconder evaporação impossível. Para sólidos, a medida validada é a massa em g; para líquidos, o volume é $massa_{consumida} / densidade_{aplicada}$, com \texttt{densidade\_aplicada} positiva e histórica do empréstimo (nunca uma densidade corrente arbitrária). Não misturar g e mL. Se \texttt{peso\_retorno} $\ge$ \texttt{peso\_saida} na devolução original, a perda bruta observada é zero e, com os dados daquela operação, a evaporação deve ser zero: não se afirma ganho líquido e evaporação positiva simultaneamente. Enquanto a pendência estiver aberta, \texttt{medida\_utilizada}, \texttt{peso\_retorno\_efetivo} permanecem \texttt{NULL} e \texttt{consumo\_validado = FALSE}; \texttt{NULL} nunca é interpretado como consumo zero. Quando uma rota estabelece um \texttt{peso\_retorno\_efetivo} diferente do original, a evaporação é \textbf{revalidada} contra a nova perda bruta antes de marcar \texttt{consumo\_validado = true}; se deixar de ser consistente, a pendência permanece (não se altera a medição nem a evaporação automaticamente, e não existe correção administrativa metrológica na V1).
\end{regra}

\begin{regra}[Tara --- três situações]
\begin{itemize}
    \item \textbf{Desconhecida} (\texttt{peso\_frasco\_vazio = NULL}, \texttt{origem\_tara = NULL}): não é zero; não estimar saldo atual apenas porque existe peso bruto.
    \item \textbf{Referência teórica} (\texttt{origem\_tara = REFERENCIA\_TEORICA}): valor derivado do nominal; não é medição física definitiva do recipiente vazio nem limite físico absoluto. Retorno inferior à referência teórica não é impossibilidade física e não deve ser ``corrigido'' artificialmente.
    \item \textbf{Real} (\texttt{origem\_tara = MEDIDA\_REAL}): só existe quando o recipiente foi efetivamente constatado vazio e pesado nessa condição. Nunca derivar de estimativa, massa nominal, densidade, produto ainda presente, cálculo reverso ou pesagem de recipiente não vazio.
\end{itemize}
\end{regra}

\begin{regra}[Esgotamento e tara real]
O esgotamento não é inferido automaticamente pela comparação com a tara; a autoridade é a confirmação explícita da condição física (\texttt{confirmarEsgotamento}), tanto para \texttt{FECHADO} quanto para \texttt{JA\_ABERTO}. Recipiente efetivamente vazio confirmado e pesado: \texttt{estado\_fisico\_frasco = VAZIO}, \texttt{disponibilidade = INDISPONIVEL}, \texttt{saldo\_desconhecido = false}, saldo corrente zero; se não havia tara real, \texttt{peso\_frasco\_vazio = peso vazio medido} e \texttt{origem\_tara = MEDIDA\_REAL}; \texttt{conteudo\_nominal} não é alterado; registra-se \texttt{FICOU\_VAZIO}. Tara real anterior divergente \textbf{não} é sobrescrita: registra-se \texttt{DISCREPANCIA\_TARA\_REAL} e a substituição exige \texttt{recalibrarTaraFrascoEsgotado} (recipiente vazio, justificativa humana, auditoria antes/depois, histórico da tara anterior preservado).
\end{regra}

\begin{regra}[Retorno abaixo da tara]
\begin{itemize}
    \item \textbf{Abaixo de tara real} (\texttt{origem\_tara = MEDIDA\_REAL}, \texttt{peso\_retorno < peso\_frasco\_vazio}, \texttt{confirmarEsgotamento = false}): preservar a leitura original; encerrar a custódia; empréstimo \texttt{DEVOLVIDO\_COM\_ANOMALIA}; pendência quantitativa aberta; frasco \texttt{INDISPONIVEL} e \texttt{em\_quarentena = true}; nunca criar saldo negativo nem tratar o retorno como zero.
    \item \textbf{Abaixo de referência teórica}: preservar a leitura; classificar \texttt{REFERENCIA\_TEORICA\_INCONSISTENTE}; não declarar impossibilidade física; seguir o contrato de inconsistência sem promover a referência a tara real.
\end{itemize}
\end{regra}

\begin{regra}[Ganho de massa --- três faixas]
A avaliação usa \texttt{peso\_retorno} contra \texttt{peso\_saida} e a tolerância Q06 determinada por \texttt{eh\_higroscopico} (snapshot imutável do frasco, preenchido no cadastro a partir de \textit{Resumo\_Reagente.eh\_higroscopico}; alterações posteriores do resumo não o reescrevem):
\begin{center}
\begin{tabular}{p{0.30\textwidth}p{0.28\textwidth}p{0.34\textwidth}}
\toprule
\textbf{Faixa} & \textbf{Sem higroscopia} & \textbf{Higroscópico} \\
\midrule
$peso_{retorno} \le peso_{saida}$ & consumo normal & consumo normal \\
$peso_{saida} < peso_{retorno} \le peso_{saida}+tol$ & consumo 0; ajuste/observação; \textbf{sem} evento de higroscopia & consumo 0; evento \texttt{ganho\_massa\_higroscopia} permitido \\
$peso_{retorno} > peso_{saida}+tol$ & \texttt{GANHO\_ACIMA\_Q06}: anomalia & \texttt{GANHO\_ACIMA\_Q06}: anomalia \\
\bottomrule
\end{tabular}
\end{center}
Ganho dentro da tolerância não produz consumo negativo e mantém o peso físico real; ganho acima de Q06 \textbf{não} provoca rollback da devolução física, encerra a custódia, gera \texttt{DEVOLVIDO\_COM\_ANOMALIA} e bloqueia o frasco (INDISPONIVEL/quarentena) para investigação. É proibido registrar evento de higroscopia para reagente não higroscópico ou com ganho fora da tolerância.
\end{regra}

\begin{regra}[Pendência metrológica --- \texttt{existePendenciaMetrologicaTx}]
O predicado é derivado de estado canônico, nunca de texto livre. Autoridade: \textit{Emprestimo\_Reagente}. Entrada: \texttt{idFrasco}. Retorna \texttt{true} se e somente se existe empréstimo daquele frasco com \texttt{status = DEVOLVIDO\_COM\_ANOMALIA} \textbf{e} \texttt{consumo\_validado = false}; caso contrário \texttt{false}. A pendência é um estado próprio do par canônico, não o status isolado. Dados fora do par (\texttt{consumo\_validado}=false com status que não seja \texttt{DEVOLVIDO\_COM\_ANOMALIA}) são \textbf{inconsistência de integridade}, não pendência metrológica: registram-se para correção e bloqueiam a liberação automática de quarentena até reconciliação.
\end{regra}

\begin{regra}[Resolução metrológica --- três rotas tipadas]
Não existe \texttt{resolverPendenciaMetrologica(campo, valor)}, dispatcher genérico público nem undo universal. A V1 possui exatamente \textbf{três rotas metrológicas tipadas} (HQ-M2-009/A, revisada por HQ-M6-001), cada uma com contrato próprio; todas preservam \texttt{peso\_retorno}, o status histórico e o vínculo \texttt{id\_resolucao\_metrologica}:
\begin{enumerate}
    \item \texttt{repetirPesagemDevolucaoAnomala}: relê frasco e empréstimo; registra a nova leitura; só valida o consumo histórico quando não houve movimentação física posterior que impeça provar o instante original (\texttt{houveMovimentacaoFisicaPosteriorTx}). Nesse caso, calcula $massa_{consumida}=(peso_{saida}-peso_{retorno}^{efetivo})-\text{Evap}$ com a densidade histórica, após revalidar a constraint de evaporação. Se houver movimentação física posterior, registra a observação e \textbf{mantém} a pendência.
    \item \texttt{confirmarEsgotamentoAposInspecao}: confirmação humana explícita e pesagem real do recipiente vazio; aplica o contrato de esgotamento; não infere vazio da anomalia; tara real divergente é preservada; o estado corrente vira \texttt{VAZIO}/INDISPONIVEL e o consumo histórico só é validado com evidência suficiente do retorno.
    \item \texttt{recalibrarTaraFrascoEsgotado}: somente com recipiente efetivamente vazio; justificativa humana; registra tara anterior e nova; não reprocessa automaticamente empréstimos históricos; não transforma a tara nova em evidência pretérita.
\end{enumerate}
\textbf{Não existe quarta rota metrológica.} A V1 \textbf{não} possui correção administrativa metrológica, nem genérica nem específica: após persistida, uma medição (\texttt{peso\_saida}, \texttt{peso\_retorno}, \texttt{peso\_retorno\_efetivo}, tara, evaporação ou consumo) não pode ser editada administrativamente, e não existem \texttt{CORRIGIR\_PESO\_RETORNO}, \texttt{ERRO\_TRANSCRICAO\_PESO}, \texttt{CORRIGIR\_EVAPORACAO}, \texttt{CORRIGIR\_TARA}, editor \texttt{campo+novoValor} ou \texttt{resolverPendenciaMetrologica} genérico. A operação \texttt{corrigirOperacao} permanece exclusivamente para tipos administrativos não metrológicos explicitamente suportados (hoje \texttt{RETIRADA\_RETIRANTE\_INCORRETO}) e \textbf{não} encerra pendência metrológica. Consequência: erros de transcrição de peso devem ser prevenidos \textbf{antes do commit} (UI-18); quando uma leitura errada é persistida e nenhuma das três rotas físicas tipadas puder legitimamente resolver a situação, a V1 \textbf{não} fabrica evidência para encerrar a pendência — a pendência permanece. Uma nova pesagem nunca é prova automática da medição histórica anterior.
A resolução grava evento \texttt{AJUSTE} vinculado ao empréstimo e à medição original, com autoria, instante, tara/densidade aplicadas e \texttt{id\_resolucao\_metrologica}. A data FLOW do consumo permanece a da devolução física; reprocessar somente as datas diretamente afetadas, sem replay nem snapshot de estoque. A resolução \textbf{não} libera a quarentena: a saída continua dependendo do contrato M5 (\texttt{resolverQuarentenaFrasco}). Cada rota é um comando idempotente no contrato global (Seção~\ref{sec:idempotencia-m7}): \texttt{idOperacao} obrigatório e identidade \texttt{(uid, tipo\_operacao, payload\_hash)}; retry devolve o resultado persistido sem reaplicar a resolução.
\end{regra}

\paragraph{Matriz de casos quantitativos.}
\begin{longtable}{p{0.28\textwidth}p{0.18\textwidth}p{0.16\textwidth}p{0.30\textwidth}}
\toprule
\textbf{Caso} & \textbf{Consumo} & \textbf{Pendência} & \textbf{Estado/quarentena} \\
\midrule
\endfirsthead
\toprule
\textbf{Caso} & \textbf{Consumo} & \textbf{Pendência} & \textbf{Estado/quarentena} \\
\midrule
\endhead
\bottomrule
\endfoot
1. retorno $<$ saída (normal) & validado, $>0$ & não & normal \\
2. retorno $=$ saída & 0 & não & normal \\
3. ganho $\le$ Q06, não higroscópico & 0 (ajuste) & não & normal \\
4. ganho $\le$ Q06, higroscópico & 0 (ajuste higroscópico) & não & normal \\
5. ganho $>$ Q06 & pendente & sim & INDISPONIVEL + quarentena \\
6. retorno $<$ referência teórica & pendente & sim & INDISPONIVEL + quarentena \\
7. retorno $<$ tara real s/ vazio & pendente & sim & INDISPONIVEL + quarentena \\
8. vazio confirmado sem tara real & validado & não & VAZIO/INDISP.; tara real nova \\
9. vazio confirmado c/ referência teórica & validado & não & VAZIO; tara real substitui teórica \\
10. vazio c/ tara real igual & validado & não & VAZIO; tara preservada \\
11. vazio c/ tara real divergente & validado & não & VAZIO; discrepância; recalibração \\
12. anomalia resolvida por pesagem válida & validado & não & status anômalo; quarentena mantida \\
13. pesagem não prova o instante histórico & pendente & sim & observação; quarentena \\
14. confirmação posterior de esgotamento & validado (se sem mov.) & não/sim & VAZIO; quarentena mantida \\
15. recalibração de tara real & não altera & conforme & antes/depois; sem replay \\
16. saldo desconhecido sem tara & pendente & conforme & saldo \texttt{NULL}; nunca zero \\
17. líquido c/ densidade histórica válida & validado em mL & não & normal \\
18. líquido sem densidade histórica válida & rejeitado & --- & operação recusada \\
19. resolução + permanência em quarentena & validado & não & INDISPONIVEL; M5 decide saída \\
20. extravio c/ \texttt{massa\_perda\_estimada\_g} & não & não & ENCERRADO\_EXTRAORD.; estimativa \\
21. evaporação $=$ perda bruta & 0 & não & consumo zero; válido \\
22. evaporação $<$ perda bruta & $>0$ & não & válido \\
23. evaporação $>$ perda bruta & rejeitado & --- & \textbf{bloqueia antes do commit} \\
24. ganho de massa ($peso_{retorno}\ge peso_{saida}$) c/ evaporação $>0$ & rejeitado & --- & \textbf{bloqueia antes do commit} \\
25. dupla digitação igual & conforme & conforme & prossegue \\
26. dupla digitação divergente & não persistido & --- & \textbf{bloqueia antes do commit} \\
\end{longtable}

\subsubsection{Idempotência (M7)}
\label{sec:idempotencia-m7}
\begin{regra}[Taxonomia: cinco conceitos distintos]
Idempotência é a propriedade de repetir a \textbf{mesma intenção lógica} sem produzir um segundo efeito lógico. Não confundir cinco mecanismos:
\begin{itemize}
    \item \textbf{Idempotência de comando} (cliente/usuário): repetir o mesmo comando autenticado devolve o resultado já persistido e não repete efeitos de domínio — chave de operação em \textit{Operacoes}.
    \item \textbf{Deduplicação de evento} (trigger/entrega at-least-once): o mesmo \texttt{eventId} não reaplica o efeito — \textit{Eventos\_Processados}.
    \item \textbf{Chave de unicidade}: impede criar dois recursos com o mesmo valor normalizado (plaqueta, matrícula, CAS) — \textit{Chaves\_Unicas}; não é idempotência de comando.
    \item \textbf{Lock concorrente}: serializa/limita operações concorrentes — \textit{Locks\_Requisicao\_Patrimonio}; não é idempotência.
    \item \textbf{Materialização reconstruível}: recalcular a mesma janela substitui a projeção de forma absoluta, sem acumular.
\end{itemize}
Idempotência \textbf{não} substitui pré-condições, autorização, locking nem controle de concorrência de domínio (I14).
\end{regra}

\begin{regra}[Identidade de uma operação]
A identidade semântica de um comando é o triplo:
\begin{center}
\texttt{(uid, tipo\_operacao, payload\_hash)}
\end{center}
onde \texttt{uid} é a identidade autenticada no servidor (nunca fornecida pelo payload), \texttt{tipo\_operacao} é a ação e \texttt{payload\_hash} é a impressão canônica do payload semanticamente relevante. Um \texttt{idOperacao} pertence a exatamente uma identidade. Todas as verificações de retry comparam os três campos: \textbf{repetir só é retry se uid, tipo\_operacao e payload\_hash coincidirem}. Verificar apenas \texttt{uid} ou apenas \texttt{uid+payload\_hash} é insuficiente e proibido (I4--I6). Um mesmo \texttt{idOperacao} com ator, ação ou payload diferente é rejeitado de forma fechada com \texttt{ALREADY\_EXISTS}/\texttt{already-exists} (I4--I6).
\end{regra}

\begin{regra}[\texttt{idOperacao}: formato, criação e reutilização]
\begin{itemize}
    \item \textbf{Formato}: string opaca gerada no cliente com CSPRNG (por exemplo UUID v4), 1 a 128 caracteres, conjunto seguro \texttt{[A-Za-z0-9\_-]}; não participa do próprio hash nem de nenhum payload.
    \item \textbf{Origem}: criada pelo cliente \textbf{antes da primeira tentativa mutável}, persistida junto da intenção local (sessão/rascunho) e \textbf{reutilizada} durante todo retry (timeout, perda de conexão, resposta desconhecida).
    \item \textbf{Novo id}: obrigatório apenas quando o usuário inicia deliberadamente uma \textbf{nova intenção}; o cliente nunca gera novo id apenas por ter perdido a resposta (isso derrotaria a idempotência).
    \item \textbf{Vida útil}: do primeiro envio até a obtenção de um resultado conclusivo (ou descarte explícito da intenção). Não reutilizar o mesmo id para outra intenção.
\end{itemize}
\end{regra}

\begin{regra}[Canonicalização determinística do payload]
O hash é calculado por um único algoritmo global, \texttt{hashPayload(tipoOperacao, payload)} = \texttt{SHA-256(tipoOperacao + "\char`\\n" + canonicalize(payload))}. A canonicalização \texttt{canonicalize} é definida como:
\begin{lstlisting}[language=TypeScript, title={Canonicalização canônica do payload (M7)}]
function canonicalize(value: unknown): string {
  if (value === null) return "null";
  const t = typeof value;
  if (t === "boolean" || t === "number" || t === "string")
    return JSON.stringify(value);
  if (Array.isArray(value))
    return "[" + value.map(canonicalize).join(",") + "]"; // ordem preservada
  if (t === "object") {
    const obj = value as Record<string, unknown>;
    const keys = Object.keys(obj)
      .filter(k => obj[k] !== undefined) // undefined e chave ausente são equivalentes
      .sort();                            // ordenação por code unit, recursiva
    return "{" + keys.map(k => JSON.stringify(k) + ":" + canonicalize(obj[k])).join(",") + "}";
  }
  throw new Error("canonicalize: tipo não suportado");
}
function hashPayload(tipoOperacao: string, payload: Record<string, unknown>): string {
  return sha256(tipoOperacao + "\n" + canonicalize(payload));
}
\end{lstlisting}
Consequências normativas: propriedades em ordem diferente produzem o mesmo hash; objetos aninhados são ordenados recursivamente; arrays \textbf{preservam a ordem} (quando a ordem é semanticamente relevante, reordenar é payload diferente); \texttt{null} e chave ausente são \textbf{distintos}; \texttt{undefined} é equivalente à ausência e não entra no hash; números/strings/booleanos seguem \texttt{JSON.stringify} (sem arredondamento, sem trim, sem uppercase, sem conversão de unidade); \texttt{idOperacao} é excluído do payload; campos derivados somente no servidor não entram no payload. A validação de justificativa humana (mínimo após trim) \textbf{não} altera o payload original usado no hash. Não usar o \emph{replacer} em forma de array do \texttt{JSON.stringify}, pois ele descarta propriedades de objetos aninhadas.
\end{regra}

\begin{regra}[Estados e resultado de \textit{Operacoes}]
\textit{Operacoes/\{idOperacao\}}: \texttt{uid}, \texttt{tipo\_operacao}, \texttt{payload\_hash}, \texttt{status}, \texttt{resultado} (resposta mínima), \texttt{criado\_em}, \texttt{atualizado\_em} — todos escritos pelo servidor.
\begin{itemize}
    \item \textbf{Comando Firestore atômico}: o documento é criado já como \texttt{CONCLUIDA}, na mesma transação do efeito de domínio; não há janela \texttt{PENDENTE}. É o caso da maioria das mutações.
    \item \textbf{Fluxo com etapa externa pós-commit} (PDF, upload, Auth, e-mail): cria-se \texttt{PENDENTE} antes/na transação, o processador idempotente executa o efeito externo e só então grava \texttt{CONCLUIDA} com o resultado; \texttt{FALHOU} é terminal para aquela tentativa.
    \item \textbf{Retry}: \texttt{CONCLUIDA} $\rightarrow$ devolve o \texttt{resultado} persistido sem reexecutar efeitos; \texttt{PENDENTE} $\rightarrow$ retoma/reconhece sem duplicar efeito; \texttt{FALHOU} $\rightarrow$ exige nova intenção ou retry explícito conforme a política da operação.
    \item \textbf{Dados inválidos}: documento \textit{Operacoes} incompatível/corrompido (campos ausentes) é tratado fail-closed; nunca se assume sucesso.
\end{itemize}
Retry exato \textbf{não} cria novo histórico, novo contador, novo empréstimo, novo frasco, nova notificação, novo timestamp lógico nem novo efeito externo (I1--I3, I10). O \texttt{timestamp} da operação é o da criação; o retry não o atualiza.
\end{regra}

\begin{regra}[Concorrência]
Duas requisições idênticas simultâneas são serializadas pela transação Firestore: a primeira confirma e cria \textit{Operacoes}; a segunda, relendo o documento confirmado, devolve o resultado sem segundo efeito (I9). Se a transação abortar antes do commit, nenhum efeito persiste. Se o callback transacional for repetido automaticamente pelo Firestore, nenhum efeito externo pode ocorrer dentro dele. Idempotência de comando não resolve concorrência entre \textbf{intenções diferentes} (dois gestores retirando o mesmo frasco, duas revogações removendo o último responsável, duas decisões de quarentena): isso continua exigindo pré-condições, leituras transacionais e locks.
\end{regra}

\begin{regra}[Efeitos externos]
É \textbf{proibido} executar efeito irreversível (e-mail, PDF, upload, Auth, chamada de rede) diretamente no callback transacional, pois o callback pode repetir. A transação persiste o fato e um \textbf{evento/outbox/operação}; um processador idempotente realiza o efeito após o commit. Para uploads, use caminho/objeto determinístico derivado de \texttt{idOperacao} (sobrescrita idempotente) e registre o resultado em \textit{Operacoes}. Falhas são recuperadas por reconciliação; ``Firestore commitou mas o efeito externo falhou'' mantém a operação \texttt{PENDENTE} e não anuncia conclusão prematura (I11).
\end{regra}

\begin{regra}[Triggers e entrega at-least-once]
\textit{Eventos\_Processados/\{eventId\}} deduplica eventos de trigger: a transação do trigger só aplica o efeito se ainda não processou aquele \texttt{eventId} e registra o processamento no mesmo commit. A linguagem correta é ``o \textbf{efeito observado} do mesmo \texttt{eventId} é transacionalmente único / deduplicado'', e \textbf{não} ``entrega exactly-once'': Cloud Functions/Eventarc pode reentregar o evento; o que o sistema garante é que a reentrega não reaplica o efeito (I12). Dois \texttt{eventId} distintos representam fatos distintos e não são deduplicados entre si.
\end{regra}

\begin{regra}[Jobs agendados]
Dois comportamentos distintos, nunca confundidos:
\begin{itemize}
    \item \textbf{Chave determinística + \texttt{create()}}: reexecução é \textbf{no-op} se a chave já existe (ex.: notificação com docId \texttt{\{id\_emprestimo\}--\{janela\}}; \texttt{BulkWriter.create} tratando \texttt{ALREADY\_EXISTS}/\texttt{code 6}). Reexecutar não cria novo fato.
    \item \textbf{Recomputação absoluta por \texttt{DATA\_ALVO}}: reexecução \textbf{recalcula e SUBSTITUI} a projeção da chave diária, \textbf{nunca incrementa} sobre o valor existente (I13). Dadas as mesmas fontes e \texttt{versao\_calculo}, o resultado é idêntico.
\end{itemize}
Jobs devem ser seguros para reexecução, execução parcial e \textit{backfill}; nenhum job acumula sobre materialização anterior.
\end{regra}

\begin{regra}[Notificações, etiquetas e materializações]
Notificações determinísticas usam docId determinístico e \texttt{create} (ou consumo transacional de evento), sem duplicar (I3, I13). \textbf{Etiquetas/PDF}: \texttt{idOperacao} é \textbf{obrigatório} (sem fallback \texttt{Date.now()}); a geração é etapa externa pós-commit com objeto de caminho determinístico por \texttt{idOperacao}, registrando \texttt{status\_geracao} \texttt{PENDENTE}$\rightarrow$\texttt{CONCLUIDA} e o resultado na operação; retry devolve o resultado armazenado. \textbf{Materializações}: reprocessar a mesma janela substitui a projeção de forma absoluta, sem dupla soma; dados de consumo validados não são somados novamente (I13).
\end{regra}

\begin{regra}[Relação com M5 e M6 e segurança]
O contrato global aplica-se às operações de M5 (extravio/reencontro/quarentena) e às três rotas metrológicas de M6 (\texttt{repetirPesagemDevolucaoAnomala}, \texttt{confirmarEsgotamentoAposInspecao}, \texttt{recalibrarTaraFrascoEsgotado}) \textbf{sem alterar a semântica de domínio} de M5/M6. \texttt{payload\_hash} é impressão de igualdade da operação, \textbf{não} é assinatura, MAC, prova de autoria, autorização nem controle de acesso; autoria é a identidade autenticada no servidor e autorização continua nas Rules/RBAC.
\end{regra}

\paragraph{Matriz de casos de idempotência (M7).}
\begin{longtable}{p{0.36\textwidth}p{0.22\textwidth}p{0.34\textwidth}}
\toprule
\textbf{Caso} & \textbf{Resultado} & \textbf{Observação} \\
\midrule
\endfirsthead
\toprule
\textbf{Caso} & \textbf{Resultado} & \textbf{Observação} \\
\midrule
\endhead
\bottomrule
\endfoot
1. mesma chave/ator/ação/payload após commit & resultado original & retry puro \\
2. mesma chave/ator, payload diferente & REJEITAR & \texttt{ALREADY\_EXISTS} \\
3. mesma chave, ator diferente & REJEITAR & \texttt{ALREADY\_EXISTS} \\
4. mesma chave, ação diferente & REJEITAR & \texttt{ALREADY\_EXISTS} \\
5. duas requisições idênticas simultâneas & um efeito & transação serializa \\
6. resposta HTTP perdida após commit & resultado original & retry com mesmo id \\
7. erro de validação antes da transação & nenhum efeito & nada persistido \\
8. erro dentro da transação antes do commit & nenhum efeito & rollback \\
9. retry automático do callback Firestore & nenhum efeito externo & outbox pós-commit \\
10. propriedades JSON em ordem diferente & mesmo hash & canonicalização \\
11. objeto aninhado em ordem diferente & mesmo hash & recursivo \\
12. array em ordem diferente & payload diferente & ordem preservada \\
13. \texttt{null} vs campo ausente & payload diferente & distintos \\
14. campo derivado só no servidor & fora do hash & server-owned \\
15. mesmo \texttt{eventId} duas vezes & efeito único & \textit{Eventos\_Processados} \\
16. dois \texttt{eventId} distintos & dois fatos & não deduplicar \\
17. job diário duas vezes & no-op/substituição & conforme tipo \\
18. materialização mesma \texttt{DATA\_ALVO} & substitui & sem dupla soma \\
19. notificação docId determinístico existente & no-op & \texttt{create}/\texttt{ALREADY\_EXISTS} \\
20. processador externo executado novamente & não repete efeito & idempotência do processador \\
21. \textit{Operacoes} incompatível/corrompido & fail-closed & nunca assumir sucesso \\
22. rota M6 repetida após resolução & resultado original & contrato global \\
23. etiquetas após resposta perdida & mesmo PDF/resultado & caminho determinístico por id \\
24. retry de correção administrativa & resultado original & contrato global \\
\end{longtable}

\subsubsection{Estoque atual, aptidão, escassez e notificações (M8)}
\label{sec:m8-estoque-escassez-notificacoes}

\begin{regra}[Estoque atual e resumos FLOW --- autoridades distintas]
\textit{Frasco\_Reagente} é a única autoridade do \textbf{estoque atual}. Não existe materialized view persistente de estoque atual e nenhuma regra pode reconstruí-lo por encadeamento diário. É proibido qualquer cálculo do tipo $estoqueAtual = resumoDeOntem + movimentacoesDeHoje$, bem como ``posição de D depende da posição de D-1''. O estoque atual é sempre calculado diretamente sobre o estado canônico corrente.

Os resumos \textit{Resumo\_Almoxarifado\_Diario} e \textit{Resumo\_Reagente\_Diario} são exclusivamente \textbf{projeções FLOW históricas}, derivadas, descartáveis e reconstruíveis a partir de \textit{Historico\_Frasco\_Reagente} e \textit{Emprestimo\_Reagente}. Não são autoridade do estoque atual e não voltam a receber campos STOCK (quantidades de frascos, massa ou volume ``no fim do dia''). Nenhuma leitura de estoque atual usa resumo diário como fonte e nenhuma correção exclusivamente histórica invalida cache de estoque atual se não alterar o estado corrente.
\end{regra}

\begin{regra}[Agregação protegida]
A cadeia de agregação é: navegador $\rightarrow$ backend protegido $\rightarrow$ \texttt{count()}/\texttt{sum()} $\rightarrow$ \textit{Frasco\_Reagente}. O navegador não baixa frascos para contá-los, não executa a autoridade da agregação e não trata o catálogo JSON como autoridade operacional. Nenhum contador cadastral (por exemplo, o vínculo de lote) substitui o estoque atual.
\end{regra}

\begin{regra}[Saldos quantitativos server-owned]
\texttt{saldo\_aferido\_g} e \texttt{saldo\_aferido\_ml} são campos derivados server-owned, recalculados na mesma transação das alterações de peso, tara ou conhecimento. Saldo desconhecido, metrologicamente inconsistente, ausente ou não mensurável é \texttt{NULL} e \textbf{nunca} é somado como zero conhecido; vazio explicitamente confirmado produz zero conhecido. A origem da tara permanece explícita e massa (g) e volume (mL) \textbf{nunca} são somados entre si.

A dashboard de estoque distingue conceitualmente: quantidade de frascos; quantidade com saldo conhecido; quantidade com saldo desconhecido; massa conhecida (g); volume conhecido (mL). Um frasco extraviado, ainda que fisicamente \texttt{ABERTO}, não integra o saldo disponível.
\end{regra}

\begin{regra}[Predicado único de aptidão: \texttt{frascoAptoParaUso}]
Existe uma única autoridade conceitual de aptidão, usada de forma mecanicamente coerente pela dashboard de estoque, pela escassez e pela retirada. Um frasco $f$ é \textbf{apto para uso} se, e somente se, todas as condições abaixo valem:
\begin{itemize}
    \item \texttt{situacao\_localizacao = LOCALIZADO};
    \item \texttt{disponibilidade = DISPONIVEL};
    \item \texttt{estado\_fisico\_frasco} $\in$ \{\texttt{FECHADO}, \texttt{ABERTO}\};
    \item \texttt{em\_quarentena = false};
    \item validade compatível: \texttt{vencido = false} \textbf{ou} \texttt{uso\_vencido\_autorizado = true};
    \item ausência de pendência metrológica impeditiva (\texttt{existePendenciaMetrologicaTx(f) = false}) e ausência de qualquer bloqueio de M5/M6.
\end{itemize}
Consequências: \texttt{ABERTO + EXTRAVIADO} continua fisicamente aberto, mas \textbf{não é apto}; \texttt{QUEBRADO}, \texttt{VAZIO}, \texttt{DESCARTADO}, \texttt{EMPRESTADO} e \texttt{em\_quarentena} não são aptos; frasco vencido sem autorização não é apto. A retirada pode acrescentar pré-condições próprias (escopo, finalidade, aceite de uso vencido), mas \textbf{não} pode empregar uma noção diferente de aptidão. Um frasco extraviado nunca conta como apto, mesmo que a disponibilidade tenha sido alterada por engano.
\end{regra}

\begin{regra}[\texttt{qtdAptos}]
Para o par formado por almoxarifado, resumo e especificação, \texttt{qtdAptos} é a contagem dos frascos de \textit{Frasco\_Reagente} cujo \texttt{id\_almoxarifado}, \texttt{id\_resumo\_reagente} e \texttt{id\_especificacao\_reagente} coincidem com o par e para os quais \texttt{frascoAptoParaUso(f)} é verdadeiro.
\texttt{qtdAptos} é \textbf{número de frascos aptos}, não massa ou volume, e é a única base do critério institucional de escassez. Não confundir com o estoque quantitativo da dashboard (que soma massa/volume conhecidos) apenas porque ambos consultam \textit{Frasco\_Reagente}.
\end{regra}

\begin{regra}[Configuração canônica de escassez: \textit{Estoques\_Configurados}]
A configuração corrente é \texttt{Almoxarifado/\{almoxId\}/Estoques\_Configurados/\{configId\}}, com chave determinística \texttt{\{id\_resumo\_reagente\}\_\{id\_especificacao\_reagente\}}. O documento contém \texttt{id\_resumo\_reagente}, \texttt{id\_especificacao\_reagente}, \texttt{id\_almoxarifado}, \texttt{qtd\_limiar\_escassez}, \texttt{ativo} e \texttt{notificacao\_ativa}. \texttt{qtd\_limiar\_escassez} é inteiro \textbf{não negativo} e mede quantidade de frascos; não é massa nem volume. A escolha entre limiar em frascos e limiar quantitativo em g/mL está decidida por \textbf{frascos} e não pode ser alterada silenciosamente.

Semântica das flags:
\begin{itemize}
    \item \texttt{ativo = false}: o par não participa da avaliação automática de escassez.
    \item \texttt{notificacao\_ativa = false}: a escassez pode continuar existindo conceitualmente, mas o job \textbf{não emite} alerta automático; silenciar não apaga a configuração, não altera o estoque e não deve ser exibido como inexistência de estoque.
    \item \texttt{Almoxarifado.ativo = false}: o job não executa a avaliação automática para o almoxarifado; operações de encerramento ordenado continuam regidas pela subseção ``Almoxarifado inativo''.
\end{itemize}
\end{regra}

\begin{regra}[Definição de escassez e fronteira]
Há escassez para o par (almoxarifado, resumo, especificação) se, e somente se:
\texttt{qtdAptos} $<$ \texttt{qtd\_limiar\_escassez}.
A comparação é \textbf{estritamente menor}. Portanto \texttt{qtdAptos = qtdLimiar} \textbf{não} é escassez e \texttt{qtdAptos > qtdLimiar} também não. Casos de fronteira: limiar $0$ nunca gera escassez (não há \texttt{qtdAptos} $<$ 0); limiar 1 gera escassez exatamente quando há 0 frascos aptos; com 3 frascos aptos e limiar 3, não há escassez. A unidade é frascos aptos; o critério não é satisfeito por saldo g/mL.
\end{regra}

\begin{regra}[Backfill legado como pré-condição operacional]
A escassez depende de \texttt{id\_resumo\_reagente} e \texttt{id\_especificacao\_reagente}. Frascos legados sem \texttt{id\_especificacao\_reagente} não podem ser ignorados em produção sem reconciliação, sob pena de falso alerta de escassez pela omissão de frascos aptos. Caso determinístico admite backfill legítimo (por exemplo, especificação determinada por lote conhecido); caso ambíguo é pendência manual de reconciliação e \textbf{nunca} se inventa uma especificação para satisfazer o schema. Especificar o job \textbf{não} significa que o backfill já foi executado nem que a cron está ativada: a cron de escassez só é operacional após a conclusão da reconciliação.
\end{regra}

\begin{regra}[Cache de dashboard não é autoridade da escassez]
A dashboard de estoque pode usar \texttt{Sistema\_Cache\_Dashboard} com validade semântica de até $30$ segundos. O job de escassez \textbf{não} decide escassez a partir de \texttt{Sistema\_Cache\_Dashboard.resultado}: a autoridade da escassez é o estado canônico corrente de \textit{Frasco\_Reagente} consultado diretamente pelo job. Cache e escassez são contratos distintos.
\end{regra}

\begin{regra}[Notificações de escassez e idempotência]
O tipo de notificação é \texttt{ESCASSEZ\_ESTOQUE}, persistido em \texttt{Usuarios/\{uid\}/Notificacoes/\{id\}}. Os campos pertinentes são \texttt{id\_destinatario}, \texttt{papel\_destinatario}, \texttt{tipo}, \texttt{entidade\_alvo}, \texttt{id\_alvo}, \texttt{quantidade}, \texttt{lida}, \texttt{emitida\_em} e \texttt{expira\_em}. Destinatários são \textbf{apenas} os gestores efetivamente vinculados ao almoxarifado (via \textit{Gestor\_Almoxarifado\_x\_Almoxarifado}); possuir o papel global de Gestor de Almoxarifado sem vínculo não faz o usuário receber o alerta. \texttt{expira\_em} é \texttt{NULL} para \texttt{ESCASSEZ\_ESTOQUE}, que é alerta operacional persistente.

A idempotência segue o contrato global (Seção~\ref{sec:idempotencia-m7}) por \textbf{docId determinístico}, equivalente a \texttt{escassez\_\{almox\}\_\{config\}\_\{dataCivil\}}, criado com \texttt{create()} (ou mecanismo equivalente). É proibido usar \texttt{Date.now()}, UUID novo por execução ou \texttt{doc()} aleatório para alertas periódicos deduplicados. Reexecutar o job no mesmo dia é \textbf{no-op} para a chave já existente; somente \texttt{ALREADY\_EXISTS} (``code 6'') é tratado como duplicidade esperada e \textbf{erros de outra natureza não podem ser engolidos}: devem propagar para não perder alertas. Configuração no dia civil seguinte pode gerar nova notificação (a \textit{dataCivil} muda).

A \textit{dataCivil} é derivada do fuso IANA \texttt{America/Sao\_Paulo}; não usar offset fixo \texttt{-03:00} nem \texttt{setHours} dependente do fuso do processo.
\end{regra}

\begin{regra}[Notificações: leitura, escrita e ``Limpar tudo'']
O cliente lê apenas as próprias notificações; a escrita é server-owned. Marcar como lida é operação protegida de servidor. ``Limpar tudo'' \textbf{marca como lida} as notificações ativas do usuário, preservando o histórico; não é \texttt{DELETE} e não pode ser implementado por exclusão arbitrária. A taxonomia e o mecanismo de idempotência de \texttt{ESCASSEZ\_ESTOQUE} são compatíveis com os demais tipos já compartilhados (devolução preventiva, atraso, frascos vencidos, vazios, quebrados, em quarentena, a serem pesados e autoatendimento).
\end{regra}

\paragraph{Matriz de invalidação do estoque atual.}
Legenda: \textbf{S} = sim; \textbf{N} = não; \textbf{S*} = sim somente quando a operação altera a contribuição corrente (saldo, disponibilidade, localização ou validade); \textbf{C} = condicional a contrato suportado. \texttt{Alm.} = escopo \texttt{ESTOQUE\_\_ALMOX} (mais o registro de geração do almoxarifado); \texttt{Alm+R} = escopo \texttt{ESTOQUE\_\_ALMOX\_\_RESUMO}. ``Orig.$\to$dest.'' identifica os almoxarifados cujas chaves são invalidadas quando há transferência.

\setlength{\tabcolsep}{2pt}
\footnotesize
\begin{longtable}{>{\raggedright\arraybackslash}p{0.25\textwidth} >{\raggedright\arraybackslash}p{0.05\textwidth} >{\raggedright\arraybackslash}p{0.08\textwidth} >{\raggedright\arraybackslash}p{0.09\textwidth} >{\raggedright\arraybackslash}p{0.14\textwidth} >{\raggedright\arraybackslash}p{0.12\textwidth} >{\raggedright\arraybackslash}p{0.12\textwidth}}
\toprule
\textbf{Operação} & \textbf{Est.} & \textbf{Alm.} & \textbf{Alm+R} & \textbf{Orig.$\to$dest.} & \textbf{FLOW} & \textbf{Notif.} \\
\midrule
\endfirsthead
\toprule
\textbf{Operação} & \textbf{Est.} & \textbf{Alm.} & \textbf{Alm+R} & \textbf{Orig.$\to$dest.} & \textbf{FLOW} & \textbf{Notif.} \\
\midrule
\endhead
\bottomrule
\endfoot
Cadastro de frasco & S & S & S & --- & \texttt{CADASTRO} & N \\
Abertura & S & S & S & --- & \texttt{ABERTURA} & N \\
Retirada & S & S & S & --- & \texttt{SAIU}, empr. & N \\
Devolução normal & S & S & S & --- & \texttt{ENTROU} & N \\
Devolução anômala & S & S & S & --- & \texttt{ENTROU}, anom. & N \\
Resolução metrológica & S* & S* & S* & --- & datas afetadas & N \\
Esgotamento & S & S & S & --- & \texttt{FICOU\_VAZIO} & N \\
Pesagem ordinária & S* & S* & S* & --- & \texttt{AJUSTE} & N \\
Evaporação & S* & S* & S* & --- & \texttt{EVAPORACAO} & N \\
Ajuste & S* & S* & S* & --- & \texttt{AJUSTE} & N \\
Recalibração de tara & S & S & S & --- & \texttt{AJUSTE} & N \\
Quebra & S & S & S & --- & \texttt{QUEBROU} & N \\
Descarte & S & S & S & --- & \texttt{FOI\_DESCARTADO} & N \\
Extravio & S & S & S & --- & \texttt{EXTRAVIO} & N \\
Reencontro & S & S & S & --- & \texttt{REENCONTRO} & N \\
Entrada em quarentena & S & S & S & --- & quarentena & N \\
Saída de quarentena & S & S & S & --- & liberação & N \\
Autorização/revogação de uso vencido & S & S & S & --- & evento & N \\
Autorização de descarte & S & S & S & --- & \texttt{PENDENTE} & N \\
Vencimento pelo job & S & S & S & --- & \texttt{VENCEU} & S \\
Transferência/associação & C & C & C & origem e destino & conforme & N \\
Correção só histórica & N & N & N & --- & reprocessa & N \\
Job de escassez & N & N & N & --- & N & S \\
\bottomrule
\end{longtable}
\normalsize

A invalidação é transacional com a mutação e \textbf{não} recalcula a agregação no commit. As notificações automáticas efetivamente especificadas na V1 para o escopo M8 são \texttt{FRASCOS\_VENCIDOS} (job de vencimento) e \texttt{ESCASSEZ\_ESTOQUE}; os demais tipos existem na taxonomia compartilhada, sem emissão automática redesenhada nesta rodada.

\subsubsection{Condição inicial, saldo desconhecido e validade desconhecida}
\begin{regra}
\texttt{condicao\_inicial\_cadastro} é um fato histórico imutável do \textit{Frasco\_Reagente}, com enum \texttt{FECHADO} ou \texttt{JA\_ABERTO}: \texttt{FECHADO} indica que o recipiente entrou fisicamente fechado no controle do LCQUI; \texttt{JA\_ABERTO} indica que já estava aberto quando passou ao controle do LCQUI. Não se confunde com a modalidade de saldo desconhecido: \texttt{saldo\_desconhecido} é propriedade metrológica atual, enquanto \texttt{condicao\_inicial\_cadastro} é fato histórico que não muda depois.

\texttt{saldo\_desconhecido} representa o desconhecimento \textbf{atual} da quantidade remanescente e não é flag histórica. Frasco cadastrado \texttt{JA\_ABERTO} inicia com \texttt{saldo\_desconhecido = true}; abertura de FECHADO preserva a flag enquanto não houver informação metrológica suficiente; ao esvaziar e medir a tara real, torna-se \texttt{false}. Somente \texttt{VAZIO} confirmado pela rota metrológica válida resolve o desconhecimento (saldo zero, flag falsa). \texttt{QUEBRADO}, \texttt{DESCARTADO}, \texttt{EXTRAVIADO} e quarentena preservam \texttt{saldo\_desconhecido}, seja verdadeiro ou falso. Estado físico e conhecimento metrológico são dimensões distintas: quebra, descarte e indisponibilidade não são medição, não fabricam saldo zero, tara ou pesagem. Resolver desconhecimento exige evidência metrológica válida; retirar do estoque não determina a quantidade física anteriormente existente. \texttt{DESCARTADO} continua terminal operacionalmente.

\texttt{abertura\_historica\_desconhecida} é flag histórica: quando verdadeira, \texttt{data\_abertura = NULL} e \texttt{estado\_fisico\_frasco} $\neq$ \texttt{FECHADO}. Não é apagada apenas porque o estado atual mudou para \texttt{VAZIO}, \texttt{QUEBRADO} ou \texttt{DESCARTADO}, nem porque a localização mudou para \texttt{EXTRAVIADO}; significa que já houve abertura e a data da primeira abertura é desconhecida. Nunca se inventa a data.

\texttt{validade\_desconhecida} é usada também quando não se pode determinar a validade efetiva exata por ausência da data histórica de abertura (\texttt{abertura\_historica\_desconhecida = true} e \texttt{validade\_apos\_aberto\_dias}). Salvo outra informação suficiente, \texttt{validade\_desconhecida = true} e \texttt{validade\_efetiva = NULL}. \texttt{validade\_desconhecida = true} não implica \texttt{vencido = false}: pode haver prova de vencimento sem data exata (por exemplo, validade fechada conhecida e já expirada). Devem ser tratadas separadamente ``conhecer a data exata de validade'' e ``ser possível concluir que já venceu''. Não se usa \texttt{vencido = true} como mecanismo de quarentena; \texttt{em\_quarentena}, \texttt{vencido}, \texttt{estado\_fisico\_frasco}, \texttt{disponibilidade} e \texttt{saldo\_desconhecido} são dimensões distintas.
\end{regra}

\subsubsection{Justificativas humanas obrigatórias (lista fechada)}
\begin{regra}
A exigência de justificativa humana com no mínimo 20 caracteres após \texttt{trim} aplica-se \textbf{somente} aos seguintes contratos enumerados:
\begin{enumerate}
    \item entrada manual em quarentena;
    \item saída ou liberação de quarentena;
    \item decisão de quarentena durante cadastro vencido;
    \item decisão de quarentena na primeira abertura;
    \item extravio;
    \item reencontro;
    \item descarte institucional;
    \item recalibração de tara real;
    \item autoatendimento (regra existente);
    \item justificativa metodológica Q04/Q14 (preservando o máximo específico existente);
    \item correção administrativa ou operacional de ação do gestor.
\end{enumerate}
Essa exigência não é regra genérica para toda string de \textit{motivo}, \textit{detalhe} ou \textit{observacao}; mensagens automáticas do servidor são isentas. Não se impõe mínimo global de 20 caracteres em \texttt{detalhe\_status} persistido. Pesagem de rotina não entra automaticamente nessa lista.
\end{regra}

\subsubsection{Limites de Geração de Relatórios}
\begin{regra}
Para proteger o sistema contra execuções exaustivas e travamentos em ambiente \textit{serverless}, os relatórios personalizados possuem limite de período de no máximo 31 dias corridos. Além disso, nenhuma data de início ou término pode ser superior ao dia atual, e o mês/ano de relatórios fechados não podem estar no futuro.
\end{regra}

\subsubsection{Almoxarifado inativo: operações bloqueadas e encerramento ordenado}
\label{sec:regras-almoxarifado-inativo}
\begin{regra}
Quando \texttt{Almoxarifado.ativo = false}, as seguintes operações são bloqueadas pelo servidor:
\begin{itemize}
    \item cadastro de novos frascos e lotes naquele almoxarifado;
    \item novas retiradas (abertura de \textit{Emprestimo\_Reagente}) vinculadas àquele almoxarifado;
    \item vinculação de novos gestores ao almoxarifado inativo.
\end{itemize}
As seguintes operações permanecem permitidas para encerramento ordenado:
\begin{itemize}
    \item registro de devoluções de empréstimos ativos (\texttt{EM\_USO} ou \texttt{ATRASADO});
    \item descarte e atualização de status de frascos já cadastrados;
    \item geração de relatórios históricos.
\end{itemize}
A reativação do almoxarifado (\texttt{ativo = true}) é possível e não exige recriação da entidade. A desativação é auditada em \textit{Registro\_de\_Auditoria} com ação \texttt{DESATIVACAO\_ALMOXARIFADO}.
\end{regra}

\subsection{Matriz formal de obrigatoriedade e dependência de campos}

\small

\small
\setlength{\tabcolsep}{5pt}
\begin{longtable}{>{\raggedright\arraybackslash}p{0.25\textwidth} >{\raggedright\arraybackslash}p{0.27\textwidth} >{\raggedright\arraybackslash}p{0.40\textwidth}}
\toprule
\textbf{Entidade} & \textbf{Condição} & \textbf{Obrigatoriedade}\\
\midrule
\endfirsthead
\toprule
\textbf{Entidade} & \textbf{Condição} & \textbf{Obrigatoriedade} (continuação)\\
\midrule
\endhead
\bottomrule
\endfoot
Resumo Reagente & sempre & nome, tipo e natureza.\\
Estoque Mínimo Almox. & sempre & id\_especificacao\_reagente, id\_almoxarifado e limiar de escassez positivo (quantidade de frascos).\\
Resumo Reagente & requer pesagem frequente & frequência em dias positiva.\\
Especificação & SOLIDO & unidade = g; densidade opcional.\\
Especificação & LIQUIDO & unidade = ml; densidade positiva obrigatória.\\
Especificação & MISTURA & pelo menos uma composição válida.\\
Frasco & condição inicial & condicao\_inicial\_cadastro = FECHADO ou JA\_ABERTO; fato histórico imutável, distinto da condição operacional (em quarentena ou não).\\
Frasco & em quarentena & dimensão operacional separada; detalhe do status (motivo) obrigatório.\\
Frasco & validade conhecida & validade\_fechado e validade\_efetiva devem ser consistentes; se já expirada, exige decisão de destino.\\
Frasco & validade desconhecida & validade\_desconhecida = true e validade\_efetiva = NULL; não implica vencido = false.\\
Frasco & validade expirada & exige destino QUARENTENA, PENDENTE\_DE\_DESCARTE ou DISPONIVEL com uso\_vencido\_autorizado.\\
Frasco & FECHADO & Peso total positivo obrigatório; conteúdo nominal positivo quando conhecido, NULL quando desconhecido. Tara de referência somente com dados suficientes; sem nominal, tara NULL e saldo desconhecido.\\
Frasco & abertura & data\_abertura e validade\_efetiva calculadas no servidor; aplicar mínimo entre validade fechada e prazo após abertura quando ambos existirem.\\
Frasco & JA\_ABERTO & peso total atual medido e conteudo\_nominal quando conhecido/leg\'ivel (\texttt{NULL} se desconhecido, nunca zero); saldo\_desconhecido = true e tara n\~ao conhecida at\'e o esgotamento.\\
Patrimônio & sempre & id\_resumo, numero\_patrimonio, estado\_conservacao, id\_local, photo\_url, nome\_responsavel\_sei e status.\\
Patrimônio & Ja\_dado\_baixa & documento comprobatório obrigatório.\\
Requisição edição & criação & motivo e ao menos um campo proposto.\\
Requisição & aprovada/rejeitada & respondida\_em, respondente e justificativa obrigatórios.\\
Convite & turma definida & e-mail normalizado + turma.\\
Convite & sem turma & e-mail normalizado + confirmação explícita do convite global.\\
Turma & criação & nome, capacidade positiva e professor responsável.\\
\bottomrule
\end{longtable}

\subsection{Matriz de combinações Multi-Role}
As únicas combinações permitidas para uma mesma conta são:
\begin{itemize}
    \item Professor + Gestor de Bens Patrimoniais;
    \item Professor + Gestor de Almoxarifado;
    \item Professor + Gestor de Bens Patrimoniais + Gestor de Almoxarifado;
    \item Aluno + Bolsista;
    \item Aluno + Bolsista + Gestor de Bens Patrimoniais;
    \item Aluno + Gestor de Almoxarifado;
    \item Aluno + Gestor de Bens Patrimoniais;
    \item Aluno + Gestor de Bens Patrimoniais + Gestor de Almoxarifado.
\end{itemize}
DP-C01: Aluno sem cargo de Bolsista pode exercer Gestor de Almoxarifado; Bolsista e Gestor de Almoxarifado são mutuamente exclusivos. Não são permitidos Chefe Geral combinado com qualquer outro papel, Aluno + Professor, nem Aluno + Bolsista + Gestor de Almoxarifado. A concessão ou remoção de papéis deve ser transacional e auditada.

\subsection{Fluxo obrigatório para mutações críticas}
Toda operação que altera dados relevantes segue: autenticar, resolver papéis, validar escopo, reler estado atual no servidor, validar transição e constraints, calcular derivados, persistir atomicamente quando necessário, registrar histórico/auditoria e disparar notificações/materializações. O frontend nunca é a autoridade final.


\subsection{Atribuição e revogação de papéis}
\label{sec:regras-role}
\subsubsection{RN-ROLE-01 — Chefe Geral não pode possuir múltiplos papéis}
O papel Chefe\_Geral é mutuamente exclusivo com todos os demais papéis do sistema.
Um usuário que possua Chefe\_Geral não pode possuir simultaneamente:
Professor;
Aluno;
Bolsista;
Gestor\_Almoxarifado;
Gestor\_Bens\_Patrimoniais;
qualquer outro papel do sistema.
Consequentemente, um usuário com Chefe\_Geral nunca é considerado Multi-Role.

\subsubsection{RN-ROLE-02 — Chefe Geral não pode revogar o papel de outro Chefe Geral}
Um usuário com o papel Chefe\_Geral não pode alterar ou revogar o papel Chefe\_Geral de outro usuário que também possua esse papel.
Um Chefe Geral somente pode iniciar o processo de revogação do próprio papel.

\subsubsection{RN-ROLE-03 — Auto-revogação de Chefe Geral}
Um Chefe Geral pode revogar o próprio papel Chefe\_Geral somente quando existir pelo menos outro usuário ativo com o papel Chefe\_Geral.
A operação deve ser rejeitada quando o usuário for o único Chefe Geral ativo do sistema.
A finalidade dessa regra é impedir que o sistema fique sem qualquer usuário com capacidade de administração global.

\subsubsection{RN-ROLE-04 — Revogação de Chefe Geral desativa a conta}
Como Chefe\_Geral é um papel exclusivo, sua revogação sempre resulta na ausência de outros papéis para o usuário.
Nesse caso:
o papel Chefe\_Geral é removido;
os Custom Claims são recalculados;
a conta do usuário é marcada como desativada (ativo = false);
o usuário permanece armazenado no sistema;
seus registros históricos e de auditoria são preservados;
a conta não deve ser excluída fisicamente.

\subsubsection{RN-ROLE-05 — Almoxarifado deve possuir pelo menos um gestor}
Todo almoxarifado ativo deve possuir obrigatoriamente pelo menos um Gestor\_Almoxarifado responsável.
Um Gestor de Almoxarifado somente poderá ter seu papel revogado se, após a operação, cada almoxarifado que ele gerencia continuar possuindo pelo menos um outro Gestor de Almoxarifado.
Portanto, a operação deve verificar individualmente todos os almoxarifados vinculados ao gestor.
Se ele for o único gestor de qualquer um deles, a revogação deverá ser rejeitada.

\subsubsection{RN-ROLE-06 — Gestores de Almoxarifado não possuem exclusividade sobre movimentações}
Não existe exigência de que o mesmo Gestor de Almoxarifado que registrou uma movimentação precise posteriormente registrar a movimentação complementar.
Quando um almoxarifado possuir dois ou mais gestores:
qualquer um dos gestores autorizados poderá registrar a entrada de um frasco;
qualquer um dos gestores autorizados poderá registrar a saída;
qualquer um dos gestores autorizados poderá registrar a devolução;
qualquer um dos gestores autorizados poderá registrar outras movimentações para as quais possua permissão.
O fato de um gestor ter registrado a saída de determinado frasco não impede que outro gestor autorizado registre posteriormente sua entrada ou devolução.
A autoria da operação deve permanecer registrada para fins de auditoria.

\subsubsection{RN-ROLE-07 — Revogação de Gestor de Almoxarifado em usuário Multi-Role}
Quando um usuário possuir Gestor\_Almoxarifado e pelo menos outro papel válido, a revogação de Gestor\_Almoxarifado não desativa sua conta.
O sistema deve:
validar as regras de revogação dos almoxarifados;
remover Gestor\_Almoxarifado dos papéis do usuário;
recalcular seus Custom Claims;
manter os demais papéis;
manter a conta ativa;
remover da interface as funcionalidades exclusivas de Gestor de Almoxarifado.

\subsubsection{RN-ROLE-08 — Revogação de Gestor de Bens Patrimoniais em usuário Multi-Role}
Quando um usuário possuir Gestor\_Bens\_Patrimoniais e pelo menos outro papel válido, a revogação de Gestor\_Bens\_Patrimoniais não desativa sua conta.
O sistema deve:
verificar a condição necessária para a revogação;
remover o papel;
recalcular os Custom Claims;
preservar os demais papéis;
manter a conta ativa;
retirar da interface as funcionalidades exclusivas do Gestor de Bens Patrimoniais.

\subsubsection{RN-ROLE-09 — Gestor de Bens Patrimoniais não pode ser o único}
A revogação de Gestor\_Bens\_Patrimoniais somente será permitida quando existir pelo menos outro usuário ativo que continue exercendo esse papel.
Caso o usuário seja o único Gestor de Bens Patrimoniais ativo, a revogação deverá ser rejeitada.
A finalidade é impedir que o sistema fique sem nenhum responsável pela gestão patrimonial.

\subsubsection{RN-ROLE-10 — Revogação do último papel desativa a conta}
Quando um usuário possuir apenas um papel ativo e esse papel for revogado, sua conta deverá ser desativada.
Isso se aplica a todos os papéis. A conta não deve ser excluída.
O usuário deverá permanecer no sistema para preservação de auditoria, histórico, autoria de operações, etc.

\subsubsection{RN-ROLE-11 — Revogação de papel não significa exclusão do usuário}
A revogação de qualquer papel deve preservar a identidade histórica do usuário.
Fluxo correto: Revogar papel -> Validar regras -> Remover papel -> Recalcular conjunto -> Atualizar Custom Claims -> Possui outro papel? (Sim = ativo, Não = desativado).

\subsubsection{RN-ROLE-12 — Usuário Multi-Role}
Um usuário será considerado Multi-Role quando possuir dois ou mais papéis simultaneamente, exceto Chefe\_Geral.
A remoção de um papel de um usuário Multi-Role deve afetar somente aquele papel específico.

\subsubsection{RN-ROLE-13 — Usuário com apenas um papel}
Se o único papel de um usuário for revogado, o usuário ficará com 0 papéis e sua conta será desativada, mas o usuário permanece armazenado para fins históricos e de auditoria.

\subsubsection{RN-ROLE-14 — Recalculo dos Custom Claims}
Os Custom Claims devem representar o conjunto efetivo de papéis do usuário após a operação. O backend deve recalcular o conjunto completo de papéis autorizados e então atualizar os Custom Claims. Se o conjunto resultante for vazio, a conta deverá ser marcada como desativada.

\subsubsection{RN-ROLE-15 — Auditoria das alterações de papéis}
Toda atribuição e toda revogação de papel devem gerar registro de auditoria contendo: usuário afetado, usuário que executou a operação, papel adicionado/removido, data/hora, resultado, justificativa e situação da conta. Revogações rejeitadas também devem ser registradas.

\subsubsection{RN-ROLE-16 — Invariante Bolsista}
A revogação do papel Aluno para um usuário é TERMINANTEMENTE BLOQUEADA pelo backend caso ele possua o papel Bolsista ativo. A revogação do papel de Bolsista deverá ocorrer de maneira prévia ou simultânea à perda do vínculo discente.

\paragraph{Execução concorrente e sincronização de identidade.}
A verificação do último chefe, do último gestor patrimonial e de cada almoxarifado deve compartilhar um documento de controle transacional; duas revogações concorrentes não podem ambas remover o último responsável. Revogar Aluno exige remover Bolsista na mesma operação ou rejeitar a solicitação. O modal explicita o conjunto resultante antes de confirmar. A gravação dos papéis e da auditoria ocorre no Firestore; a atualização do Firebase Auth é uma etapa externa, recuperável e idempotente. Uma falha nessa etapa mantém a autorização bloqueada até a reconciliação, sem anunciar conclusão prematura.

\paragraph{Q04 e Q14: contrato de aceite e autoatendimento.}
Para PESQUISA\_TCC\_POS ou ESTUDO\_DEGRADACAO\_RESIDUOS com frasco vencido ou validade desconhecida liberado com termo, obter aceite em sessão autenticada do retirante. Persistir tcr\_versao, tcr\_aceito\_em (servidor), tcr\_aceito\_por e tcr\_auditoria\_id no empréstimo, com evento correspondente em Registro\_de\_Auditoria; exigir justificativa\_metodologica de 20--2000 caracteres, verificando o mínimo após trim. Modal/checkbox são UX: backend valida versão vigente, identidade e vínculo do aceite ao frasco, finalidade e operação antes de confirmar a retirada. O TCR registra ciência e responsabilidade e não substitui regras institucionais de segurança química. Autoatendimento só quando não houver outro gestor ativo, calculado pelo servidor com justificativa humana obrigatória de pelo menos 20 caracteres após trim e notificação à chefia. Registrar abertura\_historica\_desconhecida quando a data anterior não for conhecida; nunca inventar a data.

\paragraph{Retenção V1 (DP-D02).}
Retenção indefinida de históricos e auditoria é política conservadora de rastreabilidade do LCQUI/UENF. A desativação lógica (\textit{soft delete}) é o padrão do sistema. A exclusão física de registros deve ocorrer somente quando houver fundamento normativo/operacional aplicável e ausência de referências impeditivas, não devendo a LGPD ser utilizada como justificativa genérica ou única base para deletar dados transacionais e históricos.

\paragraph{Proteção e Escopo do localStorage.}
O uso do \texttt{localStorage} para persistência de rascunhos de wizards deve respeitar rigorosamente os seguintes controles:
\begin{itemize}
    \item \textbf{Namespace e Segurança:} Utilizar o UID como namespace das chaves, embora isso não garanta isolamento contra XSS/same-origin. Tokens, PII, documentos fiscais completos e anexos jamais devem ser armazenados.
    \item \textbf{Ciclo de Vida:} Aplicar TTL (Time to Live) e verificá-lo tanto no \textit{boot} quanto na leitura. Ocorrerá limpeza compulsória após o sucesso do fluxo (salvamento no backend), no logout e em qualquer troca de identidade.
    \item \textbf{Robustez:} Utilizar \textit{schema versioning} e tratar cenários de múltiplas abas, limitando o escopo estritamente aos dados de catálogo (\textit{whitelist} de formulários autorizados).
\end{itemize}

\paragraph{Concessão e contadores de papéis (AUDITORIA-7, C008).}
Toda concessão de papel é transacional e compartilha \texttt{Controle\_Papeis/singleton} com a revogação. Ao criar papel ativo de Chefe Geral ou Gestor de Bens Patrimoniais, incrementar respectivamente \texttt{chefes\_ativos} ou \texttt{gestores\_patrimoniais\_ativos}, na mesma transação da identidade, papel e auditoria. Repetir concessão já efetivada não incrementa novamente. Inicialização e reconciliação dos contadores precedem o uso; Auth/Custom Claims são sincronizados depois do commit.

\begin{regra}[Decisões humanas HQ-M2-004 a HQ-M2-007 -- RESOLVED]
Cadastro FECHADO com conteúdo nominal desconhecido é permitido, com tara NULL e saldo desconhecido; abrir ou pesar bruto não resolve a ausência de informação. O ciclo de vida sem tara segue JA\_ABERTO em situação equivalente. Elimina-se integralmente o antigo limiar fixo de devolução/esgotamento; Q06 dinâmica permanece restrita à relação retorno/saída. Uma tara real anterior nunca é substituída silenciosamente: o esgotamento confirmado mantém saldo zero, e a discrepância segue recalibração auditável. Medição inconsistente é preservada com devolução física concluída, bloqueio operacional e consumo pendente até resolução.

Os dois resumos diários de reagentes são FLOW-only, independentes de posição anterior; correções fechadas reprocessam somente datas diretamente afetadas. Frasco\_Reagente é autoridade atual, acessada por agregações server-side protegidas e cache lazy de validade máxima 30 segundos ou invalidação anterior. Pesagem de rotina tem observação opcional, sem mínimo humano obrigatório; operações especiais conservam seus contratos. M2.2 não é iniciado por esta reconciliação documental.
\end{regra}
